\chapter{散射振幅与可观测量}\label{chap:relativistic-scattering}

外腿约化已经把具有稳定单粒子外态的相关函数转成散射振幅，实验却还要数出单位入射通量产生了多少末态。相同振幅若配上不同的连续态测度，会给出不同的量纲与数值，因此质量壳归一化、末态相空间和入射通量必须一起计算。

先用 Lorentz 不变量固定运动学，再将标量相互作用代入完整过程，能分别检验通量、全同粒子因子和单位换算。精确相关函数的谱表示随后区分稳定单粒子极点、连续阈值和传播子中的非在壳部分；圈修正的幺正性与尺度依赖由\cref{chap:quantum-radiative-structure}继续处理。

\section{散射振幅与可观测量}

有了传播子和顶角，我们仍然只有计算振幅的语言；实验并不直接测量 Green 函数。LSZ、质量壳相空间和入射通量把相关函数依次变成约化振幅、截面与衰变率，从而第一次闭合“场算符—散射振幅—实验数字”的链条。

散射振幅只有与协变态归一化、质量壳测度、Mandelstam 不变量、末态相空间和不变入射流量结合后才成为衰变率或截面。以下公式统一使用物理四动量并显式保留 $c,\hbar$，使振幅归一化与实验单位之间的关系保持可追踪。

散射积分使用 $\delta[f]$、Feynman/Schwinger 参数以及相应积分恒等式。

\subsection{相空间与 Lorentz 不变量}
从振幅到实验量的第一步由运动学决定：协变态归一化固定外态的 $\delta$ 归一化，Mandelstam 变量编码 Lorentz 不变量，两体相空间给出末态态密度，不变通量给出单位入射流的归一化。具体相互作用进入这些通用因子的唯一位置是约化矩阵元 $\mathcal M$。
% -----------------------------------------------------------------------------
% Scattering foundations in strict SI units
% -----------------------------------------------------------------------------

场论计算首先给出约化振幅 $\mathcal M$，而实验记录的是事件率和截面。两者之间的桥梁完全由运动学与归一化决定：连续末态贡献质量壳相空间，总平移不变性给出四动量守恒的 $\delta$ 分布，最后除以入射通量才得到截面。因此一旦这套 SI 制相空间公式固定，QED、QCD 与标准模型的区别只体现在代入的 $\mathcal M$。以下四动量始终采用全书统一的 SI 物理四动量
\begin{equation}
 p^\mu=\left(\frac{E_{\vb{p}}}{c},\vb{p}\right),
 \qquad
 p^2=m^2c^2,
 \qquad
 [p^\mu]=\mathrm{kg\,m\,s^{-1}}.
 \label{eq:si-scattering-fourmomentum-r9}
\end{equation}
因此 Mandelstam 变量具有动量平方量纲；截面中的面积量纲由显式的 \(\hbar^2\) 提供。四维积分限制到正能质量壳以后使用的基本测度是


\begin{proposition}[正能质量壳测度的积分化简]
设四动量为时间向量且采用正能支，$\delta$ 函数与坐标单位均按正文 SI 约定。则有
\begin{equation}
 {
 \dd^4 p\,\delta(p^2-m^2c^2)\,\theta(p^0)
 =\frac{c\,\dd^3 p}{2E_{\vb{p}}}.}
\label{eq:positive-energy-mass-shell-measure-r68}
\end{equation}
\end{proposition}
\begin{proof}
先从四维积分开始，并逐项检查其 Lorentz 变换性质：

\begin{equation*}
 I[F]\equiv
 \int\dd^4 p\,
 \delta(p^2-m^2c^2)\,\theta(p^0)F(p),
 \end{equation*}
其中 $\dd^4 p$ 的雅可比行列式在保向正时 Lorentz 变换下为 $1$，$p^2$ 是 Lorentz 标量，而这类变换保持 $p^0>0$。因此 $I[F]$ 对标量 $F$ 是 Lorentz 不变量。固定三动量后，把 $\delta$ 函数的自变量写成
\begin{equation*}
 f(p^0)=(p^0)^2-\vb{p}^2-m^2c^2.
\end{equation*}
两个根为 $p^0_\pm=\pm E_{\vb{p}}/c$，且
\begin{equation*}
 \left|\frac{\dd f}{\dd p^0}\right|_{p^0=p^0_\pm}
 =2\left|p^0_\pm\right|
 =\frac{2E_{\vb{p}}}{c}.
\end{equation*}
由一维恒等式
\begin{equation*}
 \delta[f(x)]
 =\sum_i\frac{\delta(x-x_i)}{|f'(x_i)|}
\end{equation*}
得到
\begin{equation*}
 \delta(p^2-m^2c^2)\theta(p^0)
 =\frac{c}{2E_{\vb{p}}}
 \delta\!\left(p^0-\frac{E_{\vb{p}}}{c}\right).
\end{equation*}
对 $p^0$ 积分以后，
\begin{equation*}
 I[F]
 =\int\frac{\dd^3 p}{2p^0}F(p)
 =\int\frac{c\,\dd^3 p}{2E_{\vb{p}}}F(p).
 \end{equation*}
左端是 Lorentz 不变量，因此右端就是正能质量壳上的不变量测度；完备关系~\eqnrefs{eq:rel-completeness-si-r9} 再按 Fourier 约定除以 $(2\pi\hbar)^3$。
\end{proof}

它说明连续末态的“态数”并不是普通的 \(\dd^3 p\)；因子 \(1/(2E_{\vb{p}})\) 正是 Lorentz 协变计数所要求的质量壳 Jacobian。

连续态统一采用 Wigner 单粒子表示的协变态归一化与完备关系 \eqnrefs{eq:rel-state-normalization-si-r9,eq:rel-completeness-si-r9}。Feynman 图经 LSZ 截肢后得到的约化不变量振幅记为 \(\mathcal M\)：外线 Dirac 旋量满足 \(\bar uu=2mc\)，而 \(2\to2\) 过程的 \(\mathcal M\) 取为无量纲。这个约定一旦与外腿留数一起固定，\(\hbar\) 与 \(c\) 的因子就不能再在振幅和截面之间任意移动。

对
\begin{equation*}
 p_1+p_2\longrightarrow p_3+p_4
\end{equation*}
定义
\begin{equation}
 {
 s=(p_1+p_2)^2,
 \qquad
 t=(p_1-p_3)^2,
 \qquad
 u=(p_1-p_4)^2.}
 \label{eq:mandelstam-def-si-r9}
\end{equation}
把三个定义相加并先只用四动量守恒 $p_1+p_2=p_3+p_4$，有
\begin{align*}
s+t+u
&=(p_1+p_2)^2+(p_1-p_3)^2+(p_1-p_4)^2\\
&=p_1^2+p_2^2+p_3^2+p_4^2,
\end{align*}
其中交叉项用 $p_2-p_3-p_4=-p_1$ 合并。再对四条外腿使用 $p_i^2=m_i^2c^2$，得到
\begin{equation}
 {
 s+t+u=(m_1^2+m_2^2+m_3^2+m_4^2)c^2.}
 \label{eq:stu-sum-si-r9}
\end{equation}
因此三个 Mandelstam 变量中只有两个独立；这是纯运动学约束，不依赖相互作用类型。
且质心总能量为
\begin{equation}
 E_{\rm cm}=c\sqrt{s}.
 \label{eq:Ecm-s-si-r9}
\end{equation}
因此 \(c\sqrt s\) 而非 \(\sqrt s\) 才具有能量量纲。


两体质心动量由 Källén 函数
\begin{equation*}
 \Lambda_K(a,b,c)\equiv a^2+b^2+c^2-2ab-2ac-2bc
\end{equation*}
统一给出：


\begin{proposition}[二体质量壳的 Kallen 恒等式]
设质心能量高于二体阈值，各外态在正能质量壳上，并采用正文对 Mandelstam 变量的定义。则有
\begin{equation}
 {
 |\vb{p}_i|
 =\frac{\sqrt{\Lambda_K(s,m_1^2c^2,m_2^2c^2)}}{2\sqrt s},
 \qquad
 |\vb{p}_f|
 =\frac{\sqrt{\Lambda_K(s,m_3^2c^2,m_4^2c^2)}}{2\sqrt s}.}
\label{eq:cm-momenta-kallen-si-r9}
\end{equation}
\end{proposition}
\begin{proof}
正文\eqnrefs{eq:Ecm-s-si-r9} 把质心总能量写成 $c\sqrt s$。在两体质心系中，两条粒子的三动量必定大小相同、方向相反，因此只剩一个未知的动量模。把两条质量壳条件与总能量守恒联立即可消去单粒子能量，得到正文\eqnrefs{eq:cm-momenta-kallen-si-r9}；根号何时变为实数同时给出二体阈值。

定义
\begin{equation*}
 \Lambda_K(a,b,c)=a^2+b^2+c^2-2ab-2ac-2bc.
 \end{equation*}
在初态质心系写
\begin{equation*}
 p_1^\mu=\left(\frac{E_1}{c},\vb{p}\right),
 \qquad
 p_2^\mu=\left(\frac{E_2}{c},-\vb{p}\right),
 \qquad
 c\sqrt s=E_1+E_2.
\end{equation*}
由
\begin{equation*}
 s=(p_1+p_2)^2=m_1^2c^2+m_2^2c^2+2p_1\!\cdot p_2
\end{equation*}
先得到
\begin{equation*}
 p_1\!\cdot p_2
 =\frac{s-(m_1^2+m_2^2)c^2}{2}.
\end{equation*}
另一方面
\begin{equation*}
 p_1\!\cdot p_2=\frac{E_1E_2}{c^2}+\vb{p}^2.
\end{equation*}
更直接的做法是先由 $P=p_1+p_2$ 与 $P^2=s$ 计算
\begin{align*}
 P\!\cdot p_1
 &=p_1^2+p_1\!\cdot p_2
 =\frac{s+(m_1^2-m_2^2)c^2}{2},\\
 P\!\cdot p_1
 &=\sqrt s\,\frac{E_1}{c},
\end{align*}
所以
\begin{equation*}
 E_1
 =\frac{c}{2\sqrt s}
 \left[s+(m_1^2-m_2^2)c^2\right].
\end{equation*}
再用 $E_1^2=m_1^2c^4+c^2p^2$，
\begin{align*}
 p^2
 &=\frac{1}{4s}
 \left[s+(m_1^2-m_2^2)c^2\right]^2-m_1^2c^2,\\
 &=\frac{\Lambda_K(s,m_1^2c^2,m_2^2c^2)}{4s}.
\end{align*}
因此
\begin{equation*}
 {
 |\vb{p}_i|
 =\frac{\sqrt{\Lambda_K(s,m_1^2c^2,m_2^2c^2)}}{2\sqrt s},
 \qquad
 |\vb{p}_f|
 =\frac{\sqrt{\Lambda_K(s,m_3^2c^2,m_4^2c^2)}}{2\sqrt s}.}
 \end{equation*}
末态阈值条件由根号非负给出 $c\sqrt s\ge(m_3+m_4)c^2$。
\end{proof}

根号非负同时给出末态阈值 \(c\sqrt s\ge(m_3+m_4)c^2\)。

为避免在四维 \(\delta\) 函数中混合动量与能量单位，定义能量四矢量
\begin{equation*}
 \mathcal P^\mu\equiv cp^\mu=(E,c\vb{p}).
\end{equation*}
二体 Lorentz 不变相空间于是定义为
\begin{equation}
 \dd\mathfrak P_2
 =(2\pi)^4\delta^{(4)}(\mathcal P_1+\mathcal P_2-\mathcal P_3-\mathcal P_4)
 \prod_{a=3,4}\frac{c^3\dd^3 p_a}{(2\pi)^3\,2E_a},
 \label{eq:phase2-si-start-r9}
\end{equation}
在质心系化为
\begin{equation}
 {
 \dd\mathfrak P_2
 =\frac1{16\pi^2}\frac{|\vb{p}_f|}{\sqrt s}\,\dd\Omega.}
 \label{eq:phase2-final-si-r9}
\end{equation}
这里 \(|\vb{p}_f|/\sqrt s\) 无量纲。

截面还需要一个“单位入射束流”的定义。非相对论情况下常用两束粒子的相对速度，但普通三速度差不是 Lorentz 标量。由两条入射四动量可以构造
\begin{equation}
 \mathcal I_{12}
 \equiv\sqrt{(p_1\!\cdot p_2)^2-m_1^2m_2^2c^4}
 \label{eq:invariant-flux-scalar-r9}
\end{equation}
它在任意惯性系给出同一个入射通量因子。把它换算成速度量纲后定义
\begin{equation}
 v_{\rm M}
 \equiv
 \frac{c^3\mathcal I_{12}}{E_1E_2}.
 \label{eq:moller-velocity-si-r9}
\end{equation}
这个量传统上称为 M\o ller 相对速度；名称本身并不增加新的物理假设。它在质心系满足 \(\mathcal I_{12}=|\vb{p}_i|\sqrt s\)，并恰好退化成两条轨迹的相对相遇速度。

连续态的平方 \(\delta\) 函数必须由有限空间盒和有限观测时间定义。若初态与末态的总能量分别为 $E_i$ 与 $E_f$，先定义能量失配
\begin{equation*}
 \delta E\equiv E_f-E_i,\qquad [\delta E]=\mathrm J.
\end{equation*}
能量守恒就是 $\delta E=0$。取体积 \(V\) 与观测时间 \(T\)，则
\begin{equation}
 (2\pi\hbar)^3\delta^{(3)}(\mathbf0)=V,
 \qquad
 \sum_{\vb{p}}\longrightarrow
 V\int\frac{\dd^3 p}{(2\pi\hbar)^3},
 \label{eq:box-delta-density-r15}
\end{equation}
以及
\begin{equation}
 \lim_{T\to\infty}\frac1T
 \left|\int_{-T/2}^{T/2}\dd t\,\ee^{\ii\delta E t/\hbar}\right|^2
 =2\pi\hbar\delta(\delta E).
 \label{eq:delta-square-time-r15}
\end{equation}
有限时空调节先把 $S$ 矩阵元的平方变成“单位时间、单位盒中一对入射粒子”的跃迁率；这还不是截面，因为其中仍含任意盒体积。只有再除以同一归一化下的入射通量密度，$V$ 才会消失并留下可与实验比较的面积。由此先得到二体末态跃迁率
\begin{equation}
 {
 \dd\Gamma_{1+2\to2}
 =\frac{\hbar^2c^3}{4VE_1E_2}
 \overline{|\mathcal M|^2}\,\dd\mathfrak P_2.}
 \label{eq:box-transition-rate-r15}
\end{equation}
除以入射通量密度 \(v_{\rm M}/V\)，便得到后续所有 \(2\to2\) 过程共用的截面母式
\begin{equation}
 {
 \frac{\dd\sigma}{\dd\Omega}
 =\frac{\hbar^2}{64\pi^2s}
 \frac{|\vb{p}_f|}{|\vb{p}_i|}
 \overline{|\mathcal M|^2}.}
 \label{eq:2to2-cross-section-master-si-r9}
\end{equation}
固定 \(s\) 时，\(t=(p_1-p_3)^2\) 对散射角的依赖完全由二体运动学决定。把角变量换成 \(t\)，同一截面母式可以写成
\begin{equation}
 {
 \frac{\dd\sigma}{\dd t}
 =\frac{\hbar^2\overline{|\mathcal M|^2}}
 {16\pi\Lambda_K(s,m_1^2c^2,m_2^2c^2)}.}
 \label{eq:2to2-dsigma-dt-si-r68}
\end{equation}
这个形式在比较不同交换道或直接以 Mandelstam 变量积分时更方便；它没有引入新的动力学假设，只是把\eqnrefs{eq:2to2-cross-section-master-si-r9} 的角变量重新参数化。末态若含 \(N\) 个完全相同、而相空间又以有序标签积分的粒子，还需除以 \(N!\) 消除重复计数。

对等质量弹性散射 \(m_1=m_2=m_3=m_4=m\)，在质心系取
\begin{equation*}
p_1=(E/c,\vb{p}),\quad p_2=(E/c,-\vb{p}),\quad
p_3=(E/c,\vb{p}'),\quad p_4=(E/c,-\vb{p}'),
\end{equation*}
其中 $|\vb{p}'|=|\vb{p}|\equiv p$、$\vb{p}\cdot\vb{p}'=p^2\cos\theta$。于是
\begin{equation*}
(p_1-p_3)^2=-|\vb{p}-\vb{p}'|^2,
\qquad
(p_1-p_4)^2=-|\vb{p}+\vb{p}'|^2,
\end{equation*}
再用 $E^2=m^2c^4+p^2c^2$ 代入并整理，得到
\begin{equation}
 {
 t=-2p^2(1-\cos\theta),
 \qquad
 u=-2p^2(1+\cos\theta),
 \qquad
 s=4(m^2c^2+p^2).}
 \label{eq:t-u-cm-equalmass-si-r9}
\end{equation}
具体相互作用只通过 \(\mathcal M(s,t,u)\) 进入，不改变这些运动学关系。

同一有限时间归一化给出质量为 \(M\) 的一体初态衰变母式
\begin{equation}
 {
 \dd\Gamma_{1\to n}
 =\frac{1}{2M\hbar}
 \overline{|\mathcal M|^2}\,\dd\mathfrak P_n.}
 \label{eq:master-decay-rate-si-r9}
\end{equation}
散射截面与衰变率因此共享同一套连续态归一化、质量壳测度和 LSZ 振幅约定，而不是两套互不相干的记忆公式。


三体末态第一次引入一个两体末态没有的新自由度：即使总四动量已经固定，三个末态粒子之间仍可连续重新分配内部不变质量。为把这个自由度显式化，定义
\begin{equation}
 \mathcal S_{ij}
 \equiv \bigl[c(p_i+p_j)\bigr]^2
 =c^2(p_i+p_j)^2,
 \qquad [\mathcal S_{ij}]=\mathrm{J^2}.
 \label{eq:threebody-invariant-mass-energy2-dp52}
\end{equation}
设质量为 $M$ 的母粒子衰变为 $1+2+3$，并令
$q_{12}\equiv p_1+p_2$、$\mathcal Q_{12}\equiv c q_{12}$。三体 Lorentz 不变相空间本身仍由三条单粒子质量壳测度和一个总四动量 $\delta$ 函数组成：
\begin{equation}
 {
 \dd\mathfrak P_3
 =(2\pi)^4\delta^{(4)}
 (\mathcal P-\mathcal P_1-\mathcal P_2-\mathcal P_3)
 \prod_{i=1}^{3}
 \frac{\dd^3 \vb{\mathcal P}_i}{(2\pi)^3 2E_i}.}
 \label{eq:phase3-definition-dp68}
\end{equation}
要把内部不变质量 $\mathcal S_{12}$ 显式取作积分变量，只需在\eqnrefs{eq:phase3-definition-dp68} 中插入恒等分解
\begin{equation*}
 1=\int\frac{\dd\mathcal S_{12}}{2\pi}
 \int\frac{\dd^4 \mathcal Q_{12}}{(2\pi)^4}
 (2\pi)^4\delta^{(4)}(\mathcal Q_{12}-c p_1-c p_2)
 (2\pi)\delta(\mathcal Q_{12}^2-\mathcal S_{12})\theta(\mathcal Q_{12}^0),
\end{equation*}
便可把三体问题递归拆成两个二体问题：


\begin{proposition}[多体相空间的递归分解]
设各末态为正能在壳态，在总四动量 $\delta$ 中插入中间动量恒等分解；相同粒子的组合因子单独处理。则有
\begin{equation}
 {
 \dd\mathfrak P_3(P;1,2,3)
 =\frac{\dd\mathcal S_{12}}{2\pi}\,
 \dd\mathfrak P_2(P;q_{12},3)\,
 \dd\mathfrak P_2(q_{12};1,2).}
\label{eq:phase3-recursive-dp52}
\end{equation}
\end{proposition}
\begin{proof}
正文\eqnrefs{eq:phase3-definition-dp68} 是三体末态的原始 Lorentz 不变测度。内部不变质量 $\mathcal S_{12}$ 可以作为显式积分变量。引入 $\mathcal Q_{12}=\mathcal P_1+\mathcal P_2$，并插入

\begin{align*}
1
={}&\int\dd^4 \mathcal Q_{12}\,
\delta^{(4)}(\mathcal Q_{12}-\mathcal P_1-\mathcal P_2)\\
&\times\int\dd\mathcal S_{12}\,
\delta(\mathcal Q_{12}^2-\mathcal S_{12})\theta(\mathcal Q_{12}^0).
\end{align*}
正能质量壳恒等式
\begin{equation*}
\delta(\mathcal Q_{12}^2-\mathcal S_{12})\theta(\mathcal Q_{12}^0)
=\frac{1}{2E_{12}}\delta(\mathcal Q_{12}^0-E_{12})
\end{equation*}
把 $\mathcal Q_{12}$ 的四维积分化成标准单粒子质量壳测度。重新组合总四动量 $\delta$ 函数与三条末态测度后得到
\begin{align*}
\dd\mathfrak P_3
={}&\frac{\dd\mathcal S_{12}}{2\pi}
\left[\dd\mathfrak P_2(P;q_{12},3)\right]
\left[\dd\mathfrak P_2(q_{12};1,2)\right],
\end{align*}
这就是正文\eqnrefs{eq:phase3-recursive-dp52}。中间四动量 $q_{12}=p_1+p_2$ 仅用于把同一组三体末态按 $(12)+3$ 分组；其不变量 $q_{12}^2=\mathcal S_{12}$ 是三体相空间积分变量。
\end{proof}

这个分解的物理含义很直接：先把母粒子看成衰变为“质量为 $\sqrt{\mathcal S_{12}}/c^2$ 的复合中间系统 $q_{12}$”和粒子 3，再在 $q_{12}$ 的静止系中让它衰变为粒子 1 与 2。$q_{12}=p_1+p_2$ 是对同一三体运动学作 $(12)+3$ 分组时引入的复合四动量坐标。

若母粒子无自旋，或已对其自旋取平均，则整体空间方向没有物理信息。积分掉整体转动自由度后，可直接用两项独立的不变质量平方作为 Dalitz 坐标：
\begin{equation}
 {
 \dd\mathfrak P_3
 =\frac{\dd\mathcal S_{12}\,\dd\mathcal S_{23}}
 {128\pi^3 M^2c^4}.}
 \label{eq:phase3-dalitz-dp52}
\end{equation}
因此三体相空间的非平凡内容全部进入 Dalitz 区域的边界，而不是测度中的复杂权重；若实验 Dalitz 图内部出现非均匀事件密度，它来自振幅的动力学权重而不是这一 Jacobian。要得到边界，只需在 $q_{12}$ 的静止系中观察粒子 2 与粒子 3 的夹角 $\theta^*$。此时
\begin{equation}
 {
 \mathcal S_{23}
 =m_2^2c^4+m_3^2c^4
 +2\left(E_2^*E_3^*
 -c^2|\vb{p}_2^*||\vb{p}_3^*|\cos\theta^*\right).}
 \label{eq:dalitz-angular-map-dp68}
\end{equation}
星号只表示在 $q_{12}$ 静止系取值。物理角满足 $-1\le\cos\theta^*\le1$，因此固定 $\mathcal S_{12}$ 时，$\mathcal S_{23}$ 的上下界为


\begin{proposition}[Dalitz 变量的运动学边界]
设三体末态在壳且能量为正，所取不变量与正文一致；边界由允许角度 $-1\le\cos\theta\le1$ 决定。则有
\begin{align}
 \mathcal S_{23}^{\pm}(\mathcal S_{12})
 ={}&m_2^2c^4+m_3^2c^4
 +\frac{1}{2\mathcal S_{12}}
 \Bigl[
 (\mathcal S_{12}+m_2^2c^4-m_1^2c^4)
 (M^2c^4-\mathcal S_{12}-m_3^2c^4)
 \nonumber\\
 &\qquad\qquad
 \pm
 \sqrt{\Lambda_K(\mathcal S_{12},m_1^2c^4,m_2^2c^4)}
 \sqrt{\Lambda_K(M^2c^4,\mathcal S_{12},m_3^2c^4)}
 \Bigr].
\label{eq:dalitz-boundary-dp52}
\end{align}
\end{proposition}
\begin{proof}
在 $q_{12}$ 静止系中，两次二体运动学分别给

\begin{align*}
E_2^*
&=\frac{\mathcal S_{12}+m_2^2c^4-m_1^2c^4}{2\sqrt{\mathcal S_{12}}},\\
E_3^*
&=\frac{M^2c^4-\mathcal S_{12}-m_3^2c^4}{2\sqrt{\mathcal S_{12}}},\\
c|\vb{p}_2^*|
&=\frac{\sqrt{\Lambda_K(\mathcal S_{12},m_1^2c^4,m_2^2c^4)}}{2\sqrt{\mathcal S_{12}}},\\
c|\vb{p}_3^*|
&=\frac{\sqrt{\Lambda_K(M^2c^4,\mathcal S_{12},m_3^2c^4)}}{2\sqrt{\mathcal S_{12}}}.
\end{align*}
把这些量代入正文\eqnrefs{eq:dalitz-angular-map-dp68}。由于 $\mathcal S_{23}$ 对 $\cos\theta^*$ 是线性的，固定 $\mathcal S_{12}$ 时的最大值与最小值必定出现在 $\cos\theta^*=\mp1$。逐项代入便得到正文\eqnrefs{eq:dalitz-boundary-dp52}。

两个 Källén 根号还必须同时为实数，因此
\begin{equation*}
(m_1+m_2)^2c^4
\le\mathcal S_{12}
\le(M-m_3)^2c^4.
\end{equation*}
这给出 Dalitz 区域在 $\mathcal S_{12}$ 方向上的完整阈值范围。
\end{proof}

这里同一个 K\"all\'en 函数同时控制“$1+2$ 子系统能否形成”以及“该子系统能否与粒子 3 共同来自母粒子”两个阈值。

\begin{exbox}[无质量三体末态的 Dalitz 三角形]
取 $m_1=m_2=m_3=0$。求 Dalitz 平面上的物理区域；若截肢振幅为常数，求归一化事件密度。
\textbf{解。} 三项不变质量满足
\begin{equation*}
 \mathcal S_{12}+\mathcal S_{23}+\mathcal S_{13}=M^2c^4,
 \qquad
 \mathcal S_{ij}\ge0.
\end{equation*}
因此物理区域就是
\begin{equation*}
 \mathcal S_{12}\ge0,
 \qquad
 \mathcal S_{23}\ge0,
 \qquad
 \mathcal S_{12}+\mathcal S_{23}\le M^2c^4
\end{equation*}
围成的直角三角形。除去母粒子质量尺度后，定义无量纲坐标
\begin{equation*}
 x\equiv\frac{\mathcal S_{12}}{M^2c^4},
 \qquad
 y\equiv\frac{\mathcal S_{23}}{M^2c^4}.
\end{equation*}
于是三个顶点依次为 $(0,0)$、$(1,0)$、$(0,1)$，物理区域满足 $x\ge0$、$y\ge0$、$x+y\le1$，无量纲面积恰为 $1/2$：
\begin{equation*}
 \int_0^1\dd x\int_0^{1-x}\dd y=\frac12.
\end{equation*}
\begin{center}
\begin{tikzfig}[x=4.3cm,y=4.3cm]
  \draw[lecturearrow] (0,0)--(1.12,0) node[right] {$x$};
  \draw[lecturearrow] (0,0)--(0,1.12) node[above] {$y$};
  \draw[thick] (0,0)--(1,0)--(0,1)--cycle;
  \node[below left] at (0,0) {$0$};
  \node[below] at (1,0) {$1$};
  \node[left] at (0,1) {$1$};
  \node[rotate=-45,fill=white,inner sep=1pt] at (0.55,0.55) {$x+y=1$};
  \node at (0.28,0.27) {physical region};
\end{tikzfig}
\end{center}
若截肢振幅在这个区域内近似为常数，三体相空间在 $(x,y)$ 平面上的事件密度就是常数，所以归一化后的二维密度为 $2$。任何明显的条带、边缘增强或局域峰都必须来自振幅本身的传播子、共振或角动量结构，而不是来自这个无质量三体测度的运动学边界。
\end{exbox}





\subsection{交叉与色散}
同一个截肢核在不同外态指派下产生交叉过程；解析性又把不同能区的实部、虚部和阈值连接起来。因而交叉不是“换符号口诀”，色散关系也不是独立技巧，而是同一复解析振幅的两种读取方式。
LSZ 约化把时序关联函数转化为实验散射振幅以后，同一个解析振幅还包含比单一反应道更多的信息。改变某条外腿的入射/出射解释会把散射、湮灭和对产生联系起来；真实中间态打开后出现的虚部，则通过复解析性约束阈值以下的实部。前者由交叉关系描述，后者由色散关系描述。

第一个问题的答案来自外腿的解析延拓。LSZ 振幅取自同一个截肢关联函数，只是不同物理过程把不同外腿解释为入射粒子、出射粒子或相应反粒子。改变这种解释时，四动量的符号和外线波函数必须一起改变。这一关系称为\emph{交叉}。

第二个问题来自复变量解析性。虚中间态尚不能处在真实质量壳上时，振幅可以平滑地随运动学变量变化；一旦某个真实多粒子通道打开，振幅在相应阈值处出现分支点，实轴两侧的边界值不再相同。这个差值就是割线不连续度，它的虚部由真实在壳中间态决定。若振幅在割线之外解析，Cauchy 公式便把这些边界值与其它运动学区域联系起来；为了控制高能端的积分，必要时要做有限次减法。这就是色散关系的来源。

交叉关系最清楚的记号是把四点截肢连通核统一写成“所有四动量均流入顶角”的形式
\begin{equation}
 \mathcal A(p_1,p_2,p_3,p_4),
 \qquad
 p_1+p_2+p_3+p_4=0.
 \label{eq:all-incoming-amplitude-dp9}
\end{equation}
对物理过程 $1+2\to3+4$，$p_1,p_2,p_3,p_4$ 都取正能物理四动量；出射外腿在全入射核中必须反向，因此
\begin{equation}
 \boxed{
 \mathcal M_{12\to34}
 =\mathcal A(p_1,p_2,-p_3,-p_4).}
 \label{eq:physical-from-all-incoming-dp9}
\end{equation}
若把某条外腿改解释为交叉过程中的反粒子，不能只机械地把 Mandelstam 变量换名；必须同时解析延拓该外腿四动量，并把 Dirac 外线的 $u$、$v$ 或矢量场的偏振波函数换成新物理态对应的外线因子。这个联合替换才保持同一截肢核与不同物理外态之间的对应。

\begin{note}[交叉关系的使用范围]
在逐阶微扰展开中，外腿延拓可直接作用于 Feynman 图的解析表达式。完全非微扰的交叉定理还依赖局域性、谱条件和解析延拓的全局性质。\eqnrange{eq:all-incoming-amplitude-dp9}{eq:physical-from-all-incoming-dp9} 所用的是标准局域量子场论中的外腿交叉规则，而不是对任意非局域模型的无条件恒等式。
\end{note}

转向第二个问题。固定其它运动学变量，把待研究的振幅或形状因子记作复变量 $z$ 的函数 $F(z)$。这里 $z$ 可取某个 Mandelstam 不变量，因而具有动量平方量纲。真实中间态通常从某个右手阈值 $s_R$ 开始产生割线；交叉道可以在负实轴方向产生左手割线，其最高端点记为 $s_L$。在两条割线之间选一个解析点 $z_0$ 作为减法点。

减法的目的很具体：如果 $F(z)$ 在无穷远衰减得不够快，就不能直接丢掉 Cauchy 大圆。取整数 $N\ge1$，用 $F$ 在 $z_0$ 的前 $N$ 阶 Taylor 数据构造
\begin{equation}
 \boxed{
 H_N(\zeta;z)
 \equiv
 \frac{
 F(\zeta)-\displaystyle\sum_{r=0}^{N-1}
 \frac{F^{(r)}(z_0)}{r!}(\zeta-z_0)^r
 }
 {(\zeta-z_0)^N(\zeta-z)}.}
 \label{eq:subtracted-cauchy-kernel-dp68}
\end{equation}
只要 $N$ 取得足够大，使 $H_N$ 在无穷远大圆上的积分消失，就可以把闭合轮廓收缩到实轴割线和孤立极点。于是得到 $N$ 次减法色散关系
\begin{equation}
 \boxed{
 \begin{aligned}
 F(z)
 &=\sum_{r=0}^{N-1}
 \frac{F^{(r)}(z_0)}{r!}(z-z_0)^r
 +\frac{(z-z_0)^N}{\pi}
 \int_{s_R}^{\infty}\!\dd s'\,
 \frac{\operatorname{Im}F(s'+\ii0^+)}{(s'-z_0)^N(s'-z)}
 \\
 &\quad
 +\frac{(z-z_0)^N}{\pi}
 \int_{-\infty}^{s_L}\!\dd s'\,
 \frac{\operatorname{Im}F(s'+\ii0^+)}{(s'-z_0)^N(s'-z)}
 +\text{孤立极点项}.
 \end{aligned}}
 \label{eq:general-subtracted-dispersion-dp9}
\end{equation}

\begin{proposition}[带减除的 Cauchy 表示]
设函数在切割平面解析，满足 Schwarz 实解析条件 $F(z^*)=F(z)^*$，减除点 $z_0$ 位于避开切割的实轴上，减除次数使无穷远大圆积分消失，并保证所写谱积分收敛。则有以下结论（原式 \eqnrefs{eq:general-subtracted-dispersion-dp9}）：
\begin{equation*}
 \boxed{
 \begin{aligned}
 F(z)
 &=\sum_{r=0}^{N-1}
 \frac{F^{(r)}(z_0)}{r!}(z-z_0)^r
 +\frac{(z-z_0)^N}{\pi}
 \int_{s_R}^{\infty}\!\dd s'\,
 \frac{\operatorname{Im}F(s'+\ii0^+)}{(s'-z_0)^N(s'-z)}
 \\
 &\quad
 +\frac{(z-z_0)^N}{\pi}
 \int_{-\infty}^{s_L}\!\dd s'\,
 \frac{\operatorname{Im}F(s'+\ii0^+)}{(s'-z_0)^N(s'-z)}
 +\text{孤立极点项}.
 \end{aligned}}
\end{equation*}
\end{proposition}
\begin{proof}
把轮廓收缩到两条实轴割线。实解析性给出

\begin{equation*}
 F(s'-\ii0^+)=F(s'+\ii0^+)^*,
\end{equation*}
所以割线两侧的差为
\begin{equation*}
 \operatorname{Disc}F(s')
 \equiv F(s'+\ii0^+)-F(s'-\ii0^+)
 =2\ii\operatorname{Im}F(s'+\ii0^+).
\end{equation*}
右手割线因此贡献
\begin{equation*}
 \frac{(z-z_0)^N}{\pi}
 \int_{s_R}^{\infty}\!\dd s'\,
 \frac{\operatorname{Im}F(s'+\ii0^+)}
 {(s'-z_0)^N(s'-z)},
\end{equation*}
左手割线同理给出从 $-\infty$ 到 $s_L$ 的积分。最后把预先减去的 Taylor 多项式加回，并把轮廓内未并入割线的孤立极点留作显式留数项，就得到正文\eqnrefs{eq:general-subtracted-dispersion-dp9}。

若 $z$ 逼近割线，分母 $1/(s'-z)$ 按 $z\to z+\ii0^+$ 理解。其主值部分给出实部，跨越极点产生的 $\ii\pi$ 项恢复原来的虚部，因此色散式同时编码了割线两侧的完整边界值。
\end{proof}

这里的减法多项式由 $F(z_0),F'(z_0),\ldots$ 这样的低能输入或重整化条件确定；割线积分则来自真实中间态。二者的物理来源不同，不能把减法常数误认为由幺正性自动决定。

若只有从 $s_0$ 开始的右手割线，并取一次减法 $N=1$、$z_0=0$ 且 $F(0)=0$，\eqnrefs{eq:general-subtracted-dispersion-dp9} 退化为
\begin{equation}
 \boxed{
 F(z)
 =\frac{z}{\pi}
 \int_{s_0}^{\infty}\dd s'\,
 \frac{\operatorname{Im}F(s'+\ii0^+)}{s'(s'-z)}.}
 \label{eq:once-subtracted-single-cut-dp9}
\end{equation}
例如在真空极化中，时间样区域的真实 $e^-e^+$ 产生决定 $\operatorname{Im}F$，同一色散积分据此重建类空区域的实部。

对某些由光学定理保证谱权重 $\rho(s)\equiv\operatorname{Im}F(s+\ii0^+)\ge0$ 的标量投影，\eqnrefs{eq:once-subtracted-single-cut-dp9} 还给出一个直接的符号检查。令 $z=-Q^2<0$，只要相应谱矩收敛，就有
\begin{equation}
 \boxed{
 (-1)^n\frac{\dd^n}{\dd(Q^2)^n}F(-Q^2)\ge0,
 \qquad n\ge1.}
 \label{eq:spacelike-spectral-positivity-dp68}
\end{equation}

\begin{proposition}[正谱测度的类空响应符号]
假设谱密度非负，减除处方与正文一致，并且类空积分收敛。则有以下结论（原式 \eqnrefs{eq:spacelike-spectral-positivity-dp68}）：
\begin{equation*}
 \boxed{
 (-1)^n\frac{\dd^n}{\dd(Q^2)^n}F(-Q^2)\ge0,
 \qquad n\ge1.}
\end{equation*}
\end{proposition}
\begin{proof}
这一结论只讨论谱密度 $\rho(s)\equiv\operatorname{Im}F(s+\ii0^+)\ge0$ 的投影。令 $z=-Q^2$、$Q^2>0$，正文\eqnrefs{eq:once-subtracted-single-cut-dp9} 变为
\begin{equation*}
F(-Q^2)
=-\frac{Q^2}{\pi}
\int_{s_0}^{\infty}\dd s'\,
\frac{\rho(s')}{s'(s'+Q^2)}.
\end{equation*}
利用恒等式
\begin{equation*}
\frac{Q^2}{s'(s'+Q^2)}
=\frac1{s'}-\frac1{s'+Q^2},
\end{equation*}
可以把所有 $Q^2$ 依赖集中到第二项：
\begin{equation*}
F(-Q^2)
=-\frac1\pi\int_{s_0}^{\infty}\frac{\rho(s')}{s'}\dd s'
+\frac1\pi\int_{s_0}^{\infty}\frac{\rho(s')}{s'+Q^2}\dd s'.
\end{equation*}
若相应谱矩收敛，可以逐次把 $Q^2$ 导数移入积分号。一次和二次导数分别为
\begin{align*}
\frac{\dd}{\dd Q^2}F(-Q^2)
&=-\frac1\pi\int_{s_0}^{\infty}\dd s'\,
\frac{\rho(s')}{(s'+Q^2)^2}\le0,\\
\frac{\dd^2}{\dd(Q^2)^2}F(-Q^2)
&=\frac{2}{\pi}\int_{s_0}^{\infty}\dd s'\,
\frac{\rho(s')}{(s'+Q^2)^3}\ge0.
\end{align*}
继续求导时，每次只会再产生一个负号和一个正的整数系数，因此归纳得到正文\eqnrefs{eq:spacelike-spectral-positivity-dp68}。它给出了一个很实用的自洽检查：若显式圈积分声称对应同一非负谱投影，却违反这一导数符号层级，就说明符号、减法或解析延拓中至少有一处不一致。
\end{proof}

这不是任意振幅都满足的额外公理；它只适用于当前投影的谱密度确实非负的情形。

交叉关系与色散关系连接散射振幅的两个不同层面：前者把不同外态过程识别为同一截肢关联函数在不同质量壳边界上的取值，后者用割线虚部和必要的减法常数重建解析域中的其余部分。具体理论只需要把相应粒子质量、耦合与传播子代入这些一般结构。






\subsection{标量散射实验室}
最简单的标量相互作用已经足以同时产生图、振幅、衰变率、截面与静态势，并且不引入自旋和规范自由度。它因此隔离出散射论本身的运动学、相空间与解析结构；自旋和规范代数只在更复杂的费米子与光子模型中加入。
% -----------------------------------------------------------------------------
% Scalar-only reduced momentum-space rules in strict SI units
% -----------------------------------------------------------------------------

LSZ 给出的物理场 Green 函数在标量散射中可整理为一组约化图因子。“约化”表示把每张图中重复出现的 SI 场归一化常数统一提出；四动量采用物理动量量纲，SI 归一化常数显式保留在 Green 函数与外线因子中。

质量为 $m$ 的实标量场，其动量-空间 Feynman 二点函数已由正则量子化得到
\begin{equation}
 \widetilde\Delta_F(p)
 =\frac{\ii\hbar^3c}{p^2-m^2c^2+\ii0^+}.
 \label{eq:scalar-physical-propagator-common-r15}
\end{equation}
分子 $\hbar^3c$ 来自采用的 SI 场归一化与 Fourier 约定；极点的位置 $p^2=m^2c^2$ 和 $+\ii0^+$ 边界条件与这个共同常数无关。因此定义标量约化内部线
\begin{equation}
 {
 \Delta_{\rm red}(p)
 \equiv\frac{\ii}{p^2-m^2c^2+\ii0^+}.}
 \label{eq:scalar-reduced-propagator-common-r15}
\end{equation}
它的 SI 量纲是物理动量的负二次方。这个定义不会把质量 $m$ 当成动量；分母中必须写 $m^2c^2$。


对实标量场 $\varphi$，把四次相互作用写成
\begin{equation}
 {
 \Lag_{\varphi^4}^{\rm int}
 =-\frac{g_4}{4!\,\hbar c}\varphi^4,}
 \qquad [g_4]=1.
 \label{eq:phi4-si-coupling-common-r15}
\end{equation}
由于 $[\varphi]=\sqrt{\mathrm{J/m}}$，有 $[\varphi^4]=\mathrm{J^2/m^2}$，而 $[\hbar c]=\mathrm{J\,m}$，所以右端确实具有拉格朗日密度的 SI 单位 $\mathrm{J/m^3}$。Dyson 插入与 LSZ 外腿截肢合并后，一个约化 $\varphi^4$ 顶角留下 $-\ii g_4$；这与 \eqnrefs{eq:phi4-lsz-result-r7} 的 SI 归一化完全一致。


% -----------------------------------------------------------------------------
% Full scalar-interaction laboratory in strict SI units
% -----------------------------------------------------------------------------

复标量物质场与实标量媒介场提供了不含自旋指标的最小散射模型。在这个体系中，作用量、Dyson 展开、Wick 收缩、LSZ 截肢、散射振幅、相空间和截面可以在同一套归一化下完整闭合。四动量继续采用全书的 SI 约定
\begin{equation*}
 p^\mu=(E_{\vb{p}}/c,\vb{p}),\qquad p^2=m_{\rm ch}^2c^2,
\end{equation*}
质量通过质壳条件进入能量--动量关系。复标量的质量记为 $m_{\rm ch}$，实媒介粒子的静质量记为 $m_\chi$；$m_\chi$ 的 SI 单位为 $\mathrm{kg}$。

由标量场作用量的 SI 归一化有 $[\varphi_{\rm ch}]=[\chi]=\sqrt{\mathrm{J/m}}$。因此三次单项式 $\varphi_{\rm ch}^\dagger\varphi_{\rm ch}\chi$ 的系数不能写成一个无量纲常数。定义具有\emph{动量量纲}的约化三次耦合 $g_3$，并写
\begin{equation}
 {
 \Lag_{3}
 =-\frac{g_3}{\hbar\sqrt{\hbar c}}\,\varphi_{\rm ch}^\dagger\varphi_{\rm ch}\chi,}
 \qquad [g_3]=\mathrm{kg\,m\,s^{-1}}.
 \label{eq:scalar-cubic-si-r10}
\end{equation}
量纲可逐项检查：$[\hbar\sqrt{\hbar c}]=\mathrm{kg^{3/2}m^{7/2}s^{-2}}$，所以 $[g_3/(\hbar\sqrt{\hbar c})]=\mathrm{J^{-1/2}m^{-3/2}}$；再乘 $[\varphi_{\rm ch}^\dagger\varphi_{\rm ch}\chi]=\mathrm{J^{3/2}m^{-3/2}}$，正好得到 $\mathrm{J/m^3}$。

\begin{exbox}[$\chi\to\varphi_{\rm ch}+\varphi_{\rm ch}^{\ast}$ 的两体衰变]
取母媒介粒子的四动量为 $p_\chi$，末态四动量为 $p_1,p_2$：
\begin{equation*}
 p_\chi=p_1+p_2,\qquad p_\chi^2=m_\chi^2c^2,\qquad p_1^2=p_2^2=m_{\rm ch}^2c^2.
\end{equation*}
求树级约化振幅与母粒子静止系的两体衰变率，并核对其阈值条件。
\textbf{解。} 只有一个顶角，所以约化振幅为
\begin{equation}
 \ii\mathcal M=-\ii g_3,
 \qquad |\mathcal M|^2=g_3^2.
 \label{eq:scalar-cubic-decay-mphi-r15}
\end{equation}
在母粒子静止系，
\begin{equation*}
 p_\chi^\mu=(m_\chi c,\mathbf0),\qquad
 p_1^\mu=(E/c,\vb{p}),\qquad
 p_2^\mu=(E/c,-\vb{p}).
\end{equation*}
能量守恒给 $2E=m_\chi c^2$。再用 $E^2=m_{\rm ch}^2c^4+c^2p^2$，
\begin{align*}
 p^2
 &=\frac{E^2}{c^2}-m_{\rm ch}^2c^2,\\
 &=\frac{m_\chi^2c^2}{4}-m_{\rm ch}^2c^2,
\end{align*}
所以
\begin{equation}
 {
 |\vb{p}|
 =\frac{c}{2}\sqrt{m_\chi^2-4m_{\rm ch}^2},
 \qquad m_\chi\ge 2m_{\rm ch}.}
 \label{eq:scalar-cubic-decay-p-r10}
\end{equation}
使用由 Lorentz 不变量二体相空间得到的一体衰变公式 \eqnrefs{eq:master-decay-rate-si-r9}，并代入
\begin{equation*}
 \int\dd\mathfrak P_2
 =\frac{|\vb{p}|}{4\pi m_\chi c},
\end{equation*}
得到
\begin{align}
 \Gamma_{\chi\to\varphi_{\rm ch}\varphi_{\rm ch}^{\ast}}
 &=\frac{1}{2m_\chi\hbar}g_3^2
 \frac{|\vb{p}|}{4\pi m_\chi c},\\
 &={
 \frac{g_3^2}{16\pi m_\chi\hbar}
 \sqrt{1-\frac{4m_{\rm ch}^2}{m_\chi^2}}.}
 \label{eq:scalar-cubic-decay-width-r10}
\end{align}
量纲检查为 $[g_3^2/(m_\chi\hbar)]=\mathrm{s^{-1}}$。末态粒子与反粒子可区分，所以没有 $1/2!$；若两个末态是同种实标量，相空间必须再除以 $2!$。

\medskip

\eqnrefs{eq:scalar-cubic-si-r10} 已固定三次相互作用的 SI 归一化。取 Dyson 展开的最低非零阶，把三条外腿接到同一个局域插入点，并按 LSZ 约定截去外腿传播子，就得到\eqnrefs{eq:scalar-cubic-decay-mphi-r15} 中的约化顶角振幅。这里三种场在顶角上彼此可区分，因此不存在额外的 $3!$ 组合因子。

相互作用绘景散射算符的一阶项为
\begin{align*}
\hat S_3^{(1)}
&=\frac{\ii}{\hbar c}\int\dd^4 x\,\Lag_3(x)\\
&=-\frac{\ii g_3}{\hbar^2c\sqrt{\hbar c}}
\int\dd^4 x\,
\varphi_{\rm ch}^\dagger(x)\varphi_{\rm ch}(x)\chi(x).
\end{align*}
把一条 $\varphi_{\rm ch}$、一条 $\varphi_{\rm ch}^\ast$ 和一条 $\chi$ 外腿接到该插入点，树级连通三点函数的极点部分为
\begin{align*}
\widetilde G_{3,c}^{(1)}
={}&(2\pi\hbar)^4\delta^{(4)}(p_\chi-p_1-p_2)
\left(-\frac{\ii g_3}{\hbar^2c\sqrt{\hbar c}}\right)\\
&\times
\widetilde\Delta_{\rm ch}(p_1)
\widetilde\Delta_{\rm ch}(p_2)
\widetilde\Delta_\chi(p_\chi),
\end{align*}
其中
\begin{equation*}
\widetilde\Delta_a(p)
=\frac{\ii\hbar^3c}{p^2-m_a^2c^2+\ii0^+}.
\end{equation*}
树级留数为 $Z_a=1$。对三条外腿分别乘相应逆极点因子并取质量壳极限，三个传播子全部消去，留下
\begin{equation*}
\ii\mathfrak A_3
=-\frac{\ii g_3}{\hbar^2c\sqrt{\hbar c}}.
\end{equation*}
按照约化标量振幅的定义
$\mathcal M_3\equiv\hbar^2c\sqrt{\hbar c}\,\mathfrak A_3$，于是
\begin{equation*}
\ii\mathcal M_3=-\ii g_3,
\end{equation*}
即\eqnrefs{eq:scalar-cubic-decay-mphi-r15} 的第一式。
\end{exbox}

在同一三次耦合理论中，若两种可区分物质场 $\varphi_{{\rm ch},A}$、$\varphi_{{\rm ch},B}$ 分别以 $g_{3A}$、$g_{3B}$ 耦合到同一中性媒介场 $\chi$，最低阶弹性过程
\[
A(p_1)+B(p_2)\to A(p_3)+B(p_4)
\]
由一条 $t$ 道媒介传播线连接两个三点顶角。令交换四动量 $q=p_1-p_3$、$t=q^2$，则树级约化振幅为
\begin{equation}
{
\mathcal M_t
=\frac{g_{3A}g_{3B}}{m_\chi^2c^2-t}.}
\label{eq:scalar-exchange-mphi-r15}
\end{equation}
分母直接记录媒介传播子的极点位置；$g_{3A}g_{3B}$ 与分母具有相同的动量平方量纲，因此 $\mathcal M_t$ 无量纲。它同时决定相对论散射截面和低速极限中的有效势，两者只是同一交换核在不同运动学区域的读出。

\begin{exbox}[单媒介粒子交换的角分布]
取 $m_A,m_B$ 为两种物质粒子的质量。求单媒介交换的微分截面，写出等质量质心系中的显式角分布。
\textbf{解。} 由正文\eqnrefs{eq:scalar-exchange-mphi-r15}，标量外腿没有自旋求和，故
\begin{equation}
 {
 \overline{|\mathcal M_t|^2}
 =\frac{g_{3A}^2g_{3B}^2}{(m_\chi^2c^2-t)^2}.}
 \label{eq:scalar-exchange-M2-r10}
\end{equation}
对等质量弹性散射 $m_A=m_B=m_{\rm ch}$，质心系
\begin{equation*}
 s=4(m_{\rm ch}^2c^2+p^2),\qquad
 t=-2p^2(1-\cos\theta),
\end{equation*}
而 $p_f=p_i$。代入 SI 截面主式，
\begin{equation}
 {
 \frac{\dd\sigma}{\dd\Omega}
 =\frac{\hbar^2}{64\pi^2s}
 \frac{g_{3A}^2g_{3B}^2}
 {[m_\chi^2c^2+2p^2(1-\cos\theta)]^2}.}
 \label{eq:scalar-exchange-angular-r10}
\end{equation}
这里 $\hbar^2/s$ 提供面积量纲，因此该因子必须保留在 SI 截面中。
\end{exbox}

交换振幅还可以回答一个与截面不同的问题：当两粒子运动足够慢、能量转移远小于空间动量转移时，这个相对论交换过程能否被一个定态势能描述？这个问题不能只看传播子分母的形状，因为相对论外态与 Schr\"odinger 两体态采用不同归一化。需要先把同一个低速散射截面在两种描述中匹配，再把动量空间核 Fourier 变换到坐标空间。

令
\begin{equation*}
 m_{\rm red}\equiv\frac{m_A m_B}{m_A+m_B},
 \qquad
 \kappa\equiv\frac{p}{\hbar},
\end{equation*}
分别表示两体相对运动的约化质量与波数。定态 Schr\"odinger 方程可写成 Helmholtz 形式
\begin{equation}
 (\nabla^2+\kappa^2)\psi(\vb{r})
 =\frac{2m_{\rm red}}{\hbar^2}V(\vb{r})\psi(\vb{r}),
 \label{eq:sch-helmholtz-r13}
\end{equation}
其出射 Green 核 $G_0^{(+)}$ 给出 Lippmann--Schwinger 方程
\begin{equation}
 \psi(\vb{r})=\ee^{\ii\vb{\kappa}_i\cdot\vb{r}}
 -\frac{2m_{\rm red}}{\hbar^2}\int\dd^3 r'\,
 G_0^{(+)}(\vb{r}-\vb{r}')V(\vb{r}')\psi(\vb{r}').
 \label{eq:sch-lippmann-r13}
\end{equation}
在 $|q^0|/|\vb{q}|\ll1$ 与 $|\vb{p}|/(m_{A,B}c)\ll1$ 的静态低速区间，相对论外态与 Schr\"odinger 两体态的归一化必须先统一。对可区分的弹性二体通道，在前文固定的 $S$ 矩阵、相空间与空间 Fourier 约定下，若 $\mathcal M_{\rm rel}^{\rm NR}(\vb{q})$ 表示相对论约化振幅的低速矩阵元，则通用匹配关系为
\begin{equation}
{
\widetilde V_{\rm NR}(\vb{q})
=-\frac{\hbar^3}{4m_A m_Bc}\,
\mathcal M_{\rm rel}^{\rm NR}(\vb{q}).}
\label{eq:relativistic-to-nr-potential-matching-dp75}
\end{equation}
这里 $\widetilde V_{\rm NR}$ 是对空间坐标作 Fourier 变换得到的势核，量纲为 $\mathrm{J\,m^3}$；若振幅仍含自旋指标，\eqnrefs{eq:relativistic-to-nr-potential-matching-dp75} 按固定自旋通道理解为势算符的矩阵元。这个关系只属于低速、近静态、弹性的一体势描述，不适用于显著能量转移、实粒子产生或不可忽略迟滞的区域。

对当前标量交换，\eqnrefs{eq:scalar-exchange-mphi-r15} 在该极限变成
$\mathcal M_{\rm rel}^{\rm NR}=g_{3A}g_{3B}/(\vb{q}^2+m_\chi^2c^2)$，故\eqnrefs{eq:relativistic-to-nr-potential-matching-dp75} 立即给出
\begin{equation}
 {
 \widetilde V_{\rm NR}(\vb{q})
 =-\frac{\hbar^3g_{3A}g_{3B}}{4m_A m_Bc}
 \frac{1}{\vb{q}^2+m_\chi^2c^2}.}
 \label{eq:scalar-yukawa-momentum-kernel-dp68}
\end{equation}
再按全书空间 Fourier 约定变回坐标空间，得到
\begin{equation}
 {
 V_Y(r)=-\alpha_{AB}^{(s)}\frac{\hbar c}{r}\ee^{-r/\lambda_\chi},
 \qquad
 \alpha_{AB}^{(s)}\equiv\frac{g_{3A}g_{3B}}{16\pi m_A m_Bc^2},
 \qquad
 \lambda_\chi\equiv\frac{\hbar}{m_\chi c}.}
 \label{eq:scalar-yukawa-potential-dp49}
\end{equation}
无量纲 $\alpha_{AB}^{(s)}$ 控制势强度，$\lambda_\chi$ 是媒介粒子的 Compton 长度并控制作用程。$m_\chi\to0$ 时 $\lambda_\chi\to\infty$，屏蔽因子消失；因此有质量传播子的极点位置经三维 Fourier 变换直接生成指数屏蔽与有限作用程。
同一三次顶角允许
\begin{equation*}
 \chi(p_\chi)+\varphi_{\rm ch}(p)\to\chi(p_\chi')+\varphi_{\rm ch}(p').
\end{equation*}
两条外部 $\chi$ 沿物质线有两种次序，因此内部物质线四动量分别为 $p+p_\chi$ 与 $p-p_\chi'$。约化振幅是
\begin{equation}
 {
 \mathcal M
 =g_3^2\left[
 \frac1{(p+p_\chi)^2-m_{\rm ch}^2c^2+\ii0^+}
 +\frac1{(p-p_\chi')^2-m_{\rm ch}^2c^2+\ii0^+}
 \right].}
 \label{eq:scalar-compton-analogue-si-r10}
\end{equation}
由 $p^2=p'^2=m_{\rm ch}^2c^2$、$p_\chi^2=p_\chi'^2=m_\chi^2c^2$，
\begin{align*}
 (p+p_\chi)^2-m_{\rm ch}^2c^2&=m_\chi^2c^2+2p\cdot p_\chi,\\
 (p-p_\chi')^2-m_{\rm ch}^2c^2&=m_\chi^2c^2-2p\cdot p_\chi'.
\end{align*}
这组公式的目的不是建立一个新的命名模型，而是提前看清后来 Compton 散射会反复出现的拓扑结构：两个外部玻色量子沿同一条物质线连接时，存在两种不同的时间排序或线连接顺序。两项必须先在振幅层相加，再对总振幅取模平方。真正的光子 Compton 散射还会额外包含偏振矢量、Dirac 分子和规范约束；这些结构在 QED 建立后再逐项加入。

树级图的内部动量通常被外部动量守恒完全固定；一旦内部线形成闭合回路，情况就不同了。以 $\varphi^4$ 理论四点函数的 $g_4^2$ 阶 $s$ 道泡形结构为例，两条内部线只要求总四动量为 $P=p_1+p_2$。若第一条内部线取四动量 $\ell$，第二条就是 $P-\ell$，而 $\ell$ 本身仍可连续变化，因此必须对它积分：
\begin{equation}
 \ii\mathcal M_{\rm 1L}^{(s)}
 =
 \frac{(-\ii g_4)^2}{2}
 \int_{\ell}^{(d)}
 \frac{\ii}{\ell^2-m^2c^2+\ii0^+}
 \frac{\ii}{(P-\ell)^2-m^2c^2+\ii0^+}.
 \label{eq:phi4-bubble-toy-r12}
\end{equation}
这种没有被外部动量固定、必须连续积分的内部四动量称为\emph{环动量}，相应闭合内部线结构称为\emph{环图}。这里的 $1/2$ 来自前文 Wick 收缩的组合计数。对 $d=4$，当 $|\ell|\to\infty$ 时，积分核按 $\dd^4 \ell/\ell^4$ 缩放，径向部分表现为 $\int^{\Lambda}\dd\ell/\ell$；因此高动量端出现对数发散。后面所谓正规化与重整化，正是为了先给这类高动量积分一个明确数学定义，再把可测低能参数与调节器依赖分开。



\subsection{精确关联函数与谱表示}\label{subsec:exact-correlators-spectral}

树级展开给出了相互作用振幅的最低阶近似，但完整理论的传播子并不等于有限阶图的和。更基础的问题是：不截断微扰级数时，完整二点函数必须具有怎样的方程和谱结构？稳定粒子之所以能作为散射外态，正是因为完整传播子存在孤立单粒子极点；多粒子产生则对应连续谱。因而需要直接从路径积分恒等式和完整哈密顿量的谱分解约束二点函数，而不是把极点结构当成额外假设。

路径积分首先给出相关函数之间的动力学约束。积分变量的平移不改变积分值，因此生成泛函满足一个泛函微分恒等式；继续对外源求导，就得到把不同阶相关函数连接起来的 Schwinger--Dyson 层级。

谱分解回答另一个问题：一个完整传播子究竟由哪些物理激发组成。把完整哈密顿量的完备本征态插入二点函数，稳定单粒子产生孤立极点，多粒子态产生连续谱。把这些谱权重整理成 Lorentz 不变形式，就得到 Källén--Lehmann 表示。

实标量 $\varphi^4$ 模型已经包含相互作用、完整二点函数和非平凡谱结构，同时不需要额外携带自旋或规范指标，因此足以承载这两条一般结论：积分变量不变性约束相关函数，完整态分解决定传播子的谱结构。
% Generic exact correlators and spectral representation in SI units

沿用上一小节的实标量场 $\varphi(x)$ 和四次自相互作用。这样做的好处是，传播子、顶角和 LSZ 的记号都已经出现过；现在只改变问题本身：不再把相互作用截断在某个微扰阶，而是问完整相关函数必须满足哪些恒等式。$\varphi$ 的 SI 单位仍为 $[\varphi]=\sqrt{\mathrm{J/m}}$，静质量 $m_\varphi$ 的单位为 kg，无量纲耦合仍记为 $g_4$。相应拉格朗日密度是
\begin{equation}
 \boxed{
 \Lag[\varphi]
 =\frac12\partial_\mu\varphi\,\partial^\mu\varphi
 -\frac12\frac{m_\varphi^2c^2}{\hbar^2}\varphi^2
 -\frac{g_4}{4!\,\hbar c}\varphi^4.}
 \label{eq:SD-scalar-L-SI-r12}
\end{equation}
$\Lag$ 的单位为 $\mathrm{J\,m^{-3}}$。不同阶相关函数可以统一由外源泛函微分产生。定义带源作用量
\begin{equation*}
 S[\varphi]=\frac1c\int\dd^4 x\,\Lag[\varphi].
\end{equation*}
再加入一个复数值函数 $\mathcal J(x)$，只把它当作产生场插入的数学记账变量，而不把它解释成实验电流。源作用量和归一化生成泛函分别定义为
\begin{align}
 S_{\rm src}[\varphi,\mathcal J]
 &=\frac1c\int\dd^4 x\,\mathcal J(x)\varphi(x),
 \label{eq:SD-source-action-dp76}\\
 Z[\mathcal J]
 &=\frac1{\mathcal N}\int\mathcal D\varphi\,
 \exp\!\left\{\frac{\ii}{\hbar}
 [S[\varphi]+S_{\rm src}[\varphi,\mathcal J]]\right\},
 \qquad Z[0]=1.
 \label{eq:SD-generating-functional-dp76}
\end{align}
源的量纲由 $\mathcal J\varphi$ 必须具有能量密度量纲固定为
\begin{equation*}
 [\mathcal J]=\mathrm{J^{1/2}m^{-5/2}}.
\end{equation*}
对 $Z[\mathcal J]$ 作一次泛函微分会在积分中带下一次 $\varphi(x)$ 插入，因此 $n$ 点时间有序相关函数可统一记为
\begin{equation}
 G_n(x_1,\ldots,x_n)
 \equiv
 \left.
 \left(\frac{\hbar c}{\ii}\right)^n
 \frac{\delta^nZ}
 {\delta\mathcal J(x_1)\cdots\delta\mathcal J(x_n)}
 \right|_{\mathcal J=0}.
 \label{eq:SD-correlator-source-definition-dp76}
\end{equation}
这里的 $G_n$ 是同一个量子理论在不同插入点数下的相关函数，不是彼此独立定义的动力学对象。

把经典 Euler--Lagrange 算符定义为
\begin{equation}
 \boxed{
 \mathfrak E[\varphi](x)
 \equiv c\frac{\delta S}{\delta\varphi(x)}
 =-\left(\Box+\frac{m_\varphi^2c^2}{\hbar^2}\right)\varphi(x)
 -\frac{g_4}{3!\,\hbar c}\varphi^3(x).}
 \label{eq:SD-euler-operator-SI-r16}
\end{equation}
归一化生成泛函记为 $Z[\mathcal J]$，并取 $Z[0]=1$。场空间测度在平移下不变这一条件，把“经典运动方程”提升为作用在整个生成泛函上的精确函数微分方程：
\begin{equation}
 \boxed{
 \left[
 \mathfrak E\!\left[
 \frac{\hbar c}{\ii}\frac{\delta}{\delta\mathcal J}
 \right](x)
 +\mathcal J(x)
 \right]Z[\mathcal J]=0.}
 \label{eq:SD-master-SI-r12}
\end{equation}

\begin{proposition}[泛函分部积分的主恒等式]
在保证测度与边界项可控的正则化中，源泛函可微且场变量平移的 Jacobian 为一；连续场论使用相应分布与形式积分约定。则有以下结论（原式 \eqnrefs{eq:SD-master-SI-r12}）：
\begin{equation*}
 \boxed{
 \left[
 \mathfrak E\!\left[
 \frac{\hbar c}{\ii}\frac{\delta}{\delta\mathcal J}
 \right](x)
 +\mathcal J(x)
 \right]Z[\mathcal J]=0.}
\end{equation*}
\end{proposition}
\begin{proof}
由正文\eqnrefs{eq:SD-generating-functional-dp76}，在泛函积分中作无穷小平移 $\varphi\to\varphi+\delta\varphi$。测度平移不变且场空间边界项消失时，
\begin{equation*}
 0=
 \int\mathcal D\varphi\,
 \frac{\delta}{\delta\varphi(x)}
 \exp\!\left\{
 \frac{\ii}{\hbar}(S+S_{\rm src})
 \right\}.
\end{equation*}
把导数作用到指数，
\begin{equation*}
 0=
 \int\mathcal D\varphi\,
 \frac{\ii}{\hbar}
 \left[
 \frac{\delta S}{\delta\varphi(x)}
 +\frac1c\mathcal J(x)
 \right]
 \ee^{\ii(S+S_{\rm src})/\hbar}.
\end{equation*}
正文\eqnrefs{eq:SD-euler-operator-SI-r16} 给出
\begin{equation*}
 c\frac{\delta S}{\delta\varphi(x)}
 =\mathfrak E[\varphi](x).
\end{equation*}
另一方面，由正文的源生成泛函直接求导，
\begin{equation*}
 \frac{\delta Z}{\delta\mathcal J(x)}
 =\frac{\ii}{\hbar c}
 \langle\varphi(x)\rangle_{\mathcal J}Z,
\end{equation*}
所以积分中的一次场插入可由
\begin{equation*}
 \varphi(x)\longrightarrow
 \frac{\hbar c}{\ii}\frac{\delta}{\delta\mathcal J(x)}
\end{equation*}
替换。将同一替换用于 $\mathfrak E[\varphi]$ 中的各次幂，便有
\begin{equation*}
 \left[
 \mathfrak E\!\left[
 \frac{\hbar c}{\ii}\frac{\delta}{\delta\mathcal J}
 \right](x)
 +\mathcal J(x)
 \right]Z[\mathcal J]=0,
\end{equation*}
即正文\eqnrefs{eq:SD-master-SI-r12}。
\end{proof}

这就是 Schwinger--Dyson 母方程。它的物理内容是：经典场方程中的每个 $\varphi$ 因子，在量子理论里都变成对外源的泛函微分，所以一个固定点数的相关函数会与更高点数相关函数耦合。对母方程再作一次源微分并令 $\mathcal J=0$，第一层方程如下。
\begin{proposition}[四次标量相互作用的相关函数层级]
设正则化后的源泛函满足 Schwinger--Dyson 母方程、$Z[0]=1$，且源微分可交换。取正文的四次相互作用与源归一化；$G_n$ 表示完整时序相关函数，重合点插入按同一正则化定义。则一次源微分给出
\begin{equation}
 \boxed{
 \left(\Box_x+\frac{m_\varphi^2c^2}{\hbar^2}\right)G_2(x,y)
 +\frac{g_4}{3!\,\hbar c}G_4(x,x,x,y)
 =-\ii\hbar c\,\delta^{(4)}(x-y).}
 \label{eq:SD-two-four-hierarchy-dp76}
\end{equation}
\end{proposition}
\begin{proof}
母方程中“把场换成源导数”很紧凑，但第一次接触 Schwinger--Dyson 层级时，最容易混淆的是三件事：每次源微分带来什么 $\hbar,c,\ii$ 因子，$\varphi^3(x)$ 为什么产生四点函数，以及右端的 $\delta^{(4)}(x-y)$ 从哪里出现。为把这三件事分开，定义场插入算符
\begin{equation*}
\widehat\Phi_x
\equiv\frac{\hbar c}{\ii}\frac{\delta}{\delta\mathcal J(x)}.
\end{equation*}
正文\eqnrefs{eq:SD-correlator-source-definition-dp76} 因而等价于
\begin{equation*}
G_n(x_1,\ldots,x_n)
=\left.
\widehat\Phi_{x_1}\cdots\widehat\Phi_{x_n}Z[\mathcal J]
\right|_{\mathcal J=0}.
\end{equation*}
特别地，$Z[0]=1$ 给
\begin{equation*}
\left.\widehat\Phi_x\widehat\Phi_yZ\right|_0=G_2(x,y),
\qquad
\left.\widehat\Phi_x^3\widehat\Phi_yZ\right|_0=G_4(x,x,x,y).
\end{equation*}
第二式没有额外的 $3!$：正文 Euler--Lagrange 算符里的 $1/3!$ 已经来自对 $\varphi^4/4!$ 求一次场导数；这里的三次 $\widehat\Phi_x$ 只是把三个位于同一点 $x$ 的场插入逐个生成出来。

把正文\eqnrefs{eq:SD-euler-operator-SI-r16} 代入母方程，先把算符写开：
\begin{equation*}
\left[
-\left(\Box_x+\frac{m_\varphi^2c^2}{\hbar^2}\right)\widehat\Phi_x
-\frac{g_4}{3!\,\hbar c}\widehat\Phi_x^3
+\mathcal J(x)
\right]Z[\mathcal J]=0.
\end{equation*}
现在从左边再作用一次 $\widehat\Phi_y$。前两个项中的源导数彼此可交换，所以
\begin{align*}
&-\left(\Box_x+\frac{m_\varphi^2c^2}{\hbar^2}\right)
\widehat\Phi_x\widehat\Phi_yZ
-\frac{g_4}{3!\,\hbar c}
\widehat\Phi_x^3\widehat\Phi_yZ\\
&\qquad
+\widehat\Phi_y\!\left[\mathcal J(x)Z\right]=0.
\end{align*}
真正产生接触项的是最后一个显式源因子。乘积法则给
\begin{align*}
\widehat\Phi_y\!\left[\mathcal J(x)Z\right]
&=\frac{\hbar c}{\ii}
\frac{\delta\mathcal J(x)}{\delta\mathcal J(y)}Z
+\mathcal J(x)\widehat\Phi_yZ\\
&=\frac{\hbar c}{\ii}\delta^{(4)}(x-y)Z
+\mathcal J(x)\widehat\Phi_yZ.
\end{align*}
令 $\mathcal J=0$ 后第二项消失，而 $Z[0]=1$，因此
\begin{equation*}
\left.\widehat\Phi_y[\mathcal J(x)Z]\right|_0
=-\ii\hbar c\,\delta^{(4)}(x-y).
\end{equation*}
同时利用上面的相关函数定义，母方程的一次源微分成为
\begin{equation*}
-\left(\Box_x+\frac{m_\varphi^2c^2}{\hbar^2}\right)G_2(x,y)
-\frac{g_4}{3!\,\hbar c}G_4(x,x,x,y)
-\ii\hbar c\,\delta^{(4)}(x-y)=0.
\end{equation*}
整体乘以 $-1$ 就得到正文\eqnrefs{eq:SD-two-four-hierarchy-dp76}：
\begin{equation*}
\left(\Box_x+\frac{m_\varphi^2c^2}{\hbar^2}\right)G_2(x,y)
+\frac{g_4}{3!\,\hbar c}G_4(x,x,x,y)
=-\ii\hbar c\,\delta^{(4)}(x-y).
\end{equation*}
因此层级不闭合的来源也已经显式可见：四次相互作用的经典运动方程含 $\varphi^3$，对二点函数方程而言它必然生成三个 $x$ 插入加一个 $y$ 插入的 $G_4$；接触 $\delta$ 则来自对外源本身的泛函微分，而不是相互作用顶角。
\end{proof}
因此相互作用理论的二点函数通常不能独立闭合；继续微分会依次得到 $G_4$、$G_6$ 等更高相关函数的方程，形成无限层级。

对二点函数本身，另一种精确组织方式是把所有单粒子不可约二点插入重求和。记自由约化传播子为 $G_0(p)$、完整约化传播子为 $G(p)$。定义标量自能 $\Sigma_\varphi(p)$ 为所有不能通过切断一条内部标量线而分离的二点单粒子不可约插入之和。由于 $[G_0]=[p]^{-2}$，有 $[\Sigma_\varphi]=[p]^2$。按 1PI 自能插入次数组织完整二点函数，先得到形式级数
\begin{equation}
 \boxed{
 G
 =G_0+G_0\Sigma_\varphi G_0
 +G_0\Sigma_\varphi G_0\Sigma_\varphi G_0+\cdots .}
 \label{eq:exact-dyson-series-scalar-r68}
\end{equation}
把这一级数重求和后，等价的逆传播子形式为
\begin{equation}
 \boxed{G^{-1}(p)=G_0^{-1}(p)-\Sigma_\varphi(p).}
 \label{eq:exact-dyson-scalar-SI-r15}
\end{equation}

\begin{proposition}[自能插入的预解式重求和]
假设算子乘积及逆在共同定义域上有意义；级数求和先在收敛区域完成，再按相同解析处方延拓。则有以下结论（原式 \eqnrefs{eq:exact-dyson-scalar-SI-r15}）：
\begin{equation*}
 \boxed{G^{-1}(p)=G_0^{-1}(p)-\Sigma_\varphi(p).}
\end{equation*}
\end{proposition}
\begin{proof}
按自能插入次数排列完整标量二点函数。每个 $\Sigma_\varphi$ 两侧都由自由传播子连接，因此

\begin{align*}
 G
 &=G_0+G_0\Sigma_\varphi G_0
 +G_0\Sigma_\varphi G_0\Sigma_\varphi G_0+\cdots,\\
 &=G_0\left[
 1+\Sigma_\varphi G_0+(\Sigma_\varphi G_0)^2+\cdots
 \right].
\end{align*}
在几何级数作为形式微扰级数或相应逆算符存在的范围内，括号内可写成
\begin{equation*}
 \left(1-\Sigma_\varphi G_0\right)^{-1},
\end{equation*}
所以
\begin{equation*}
 G=G_0(1-\Sigma_\varphi G_0)^{-1}.
\end{equation*}
对标量动量空间核，$G_0(p)$ 与 $\Sigma_\varphi(p)$ 是同一固定四动量处的复数函数，因而彼此可交换。先取逆：
\begin{align*}
 G^{-1}
 &=\left[G_0(1-\Sigma_\varphi G_0)^{-1}\right]^{-1},\\
 &=(1-\Sigma_\varphi G_0)G_0^{-1},\\
 &=G_0^{-1}-\Sigma_\varphi G_0G_0^{-1},\\
 &=G_0^{-1}-\Sigma_\varphi.
\end{align*}
这就是\eqnrefs{eq:exact-dyson-scalar-SI-r15}。矩阵值或张量值传播子仍满足同一算符恒等式，但必须保留因子次序，不能使用标量交换性简化中间步骤。
\end{proof}

Dirac 场和规范场的对应关系分别提升为自旋矩阵和 Lorentz 张量逆核；标量情形给出这一逆核结构的最小代数原型。

Dyson 方程告诉我们“相互作用怎样改写逆传播子”，但还没有说明完整传播子允许包含哪些谱成分。为此取精确真空 $|\Omega\rangle$，并令 $|n\rangle$ 是完整 Hamilton 算符和四动量算符的共同谱态。对厄米标量场定义 Wightman 二点函数
\begin{equation*}
 W(x)\equiv\langle\Omega|\hat\varphi(x)\hat\varphi(0)|\Omega\rangle.
\end{equation*}
在两个场之间插入完整态的单位分解，并使用平移协变性，得到离散记号下的精确谱和
\begin{equation}
 \boxed{
 W(x)
 =\sum_n \ee^{-\ii p_n\cdot x/\hbar}
 |\langle\Omega|\hat\varphi(0)|n\rangle|^2.}
 \label{eq:W-spectral-sum-SI-r12}
\end{equation}
连续谱时求和符号同时代表相应 Lorentz 不变积分。矩阵元以模方出现，所以谱权重非负。把具有同一不变量
\begin{equation*}
 \zeta\equiv p_n^2,\qquad [\zeta]=[p]^2,
\end{equation*}
的态合并，定义 Källén--Lehmann 谱密度 $\rho_{\rm KL}(\zeta)$。采用全书的物理四动量归一化，Wightman 函数可写成
\begin{equation}
 \boxed{
 W(x)=\int_0^\infty\dd\zeta\,\rho_{\rm KL}(\zeta)
 \int\frac{\dd^3 p}{(2\pi\hbar)^3}\frac{\hbar^2 c}{2p^0}
 \ee^{-\ii p\cdot x/\hbar},
 \qquad p^0=\sqrt{\mathbf p^2+\zeta}.}
 \label{eq:KL-Wightman-mass-shell-r15}
\end{equation}
这里 $p^0$ 具有动量量纲，而不是能量量纲。谱密度满足
\begin{equation}
 \boxed{
 \rho_{\rm KL}(\zeta)\ge0,
 \qquad
 \int_0^\infty\dd\zeta\,\rho_{\rm KL}(\zeta)=1,
 \qquad [\rho_{\rm KL}]=[p]^{-2}.}
 \label{eq:KL-positive-normalization-r16}
\end{equation}

\begin{proposition}[谱测度的正性与归一化]
设 Hilbert 内积正定，场厄米，真空平移不变且满足谱条件；归一化使用正文规定的正则场假设。则有以下结论（原式 \eqnrefs{eq:KL-positive-normalization-r16}）：
\begin{equation*}
 \boxed{
 \rho_{\rm KL}(\zeta)\ge0,
 \qquad
 \int_0^\infty\dd\zeta\,\rho_{\rm KL}(\zeta)=1,
 \qquad [\rho_{\rm KL}]=[p]^{-2}.}
\end{equation*}
\end{proposition}
\begin{proof}
令 $r\equiv\mathbf x-\mathbf y$、
$\tau\equiv x^0-y^0$。对固定 $\zeta$，正文质量壳核为
\begin{equation*}
 W_\zeta(\tau,\mathbf r)
 =\int\frac{\dd^3p}{(2\pi\hbar)^3}
 \frac{\hbar^2c}{2p^0}
 \exp\!\left[-\frac{\ii}{\hbar}
 (p^0\tau-\mathbf p\cdot\mathbf r)\right],
 \qquad p^0=\sqrt{\mathbf p^2+\zeta}.
\end{equation*}
场对易函数中这一质量壳的贡献是
\[
 C_\zeta(\tau,\mathbf r)
 =W_\zeta(\tau,\mathbf r)-W_\zeta(-\tau,-\mathbf r).
\]
现在对 $x^0$ 求导；因为 $\partial_{x^0}\tau=1$，第一项给
\begin{align*}
 \left.\partial_{x^0}W_\zeta(\tau,\mathbf r)\right|_{\tau=0}
 &=-\frac{\ii\hbar c}{2}
 \int\frac{\dd^3p}{(2\pi\hbar)^3}
 \ee^{+\ii\mathbf p\cdot\mathbf r/\hbar}.
\end{align*}
第二项要同时注意前面的减号和 $-\tau$ 的链式法则：
\begin{align*}
 \left.\partial_{x^0}[-W_\zeta(-\tau,-\mathbf r)]\right|_{\tau=0}
 &=-\frac{\ii\hbar c}{2}
 \int\frac{\dd^3p}{(2\pi\hbar)^3}
 \ee^{-\ii\mathbf p\cdot\mathbf r/\hbar}.
\end{align*}
在第二个积分中令 $\mathbf p\mapsto-\mathbf p$，两项相同，因此
\begin{align*}
 \left.\partial_{x^0}C_\zeta(\tau,\mathbf r)\right|_{\tau=0}
 &=-\ii\hbar c
 \int\frac{\dd^3p}{(2\pi\hbar)^3}
 \ee^{+\ii\mathbf p\cdot\mathbf r/\hbar}\\
 &=-\ii\hbar c\,\delta^{(3)}(\mathbf r).
\end{align*}
注意 $p^0$ 已在求导时与核中的 $1/(2p^0)$ 精确约掉，所以这个等时结果与 $\zeta$ 无关。把所有质量壳按谱密度叠加：
\begin{align*}
 \left.\partial_{x^0}C(x-y)\right|_{x^0=y^0}
 =-\ii\hbar c\,
 \left[\int_0^\infty\dd\zeta\,\rho_{\rm KL}(\zeta)\right]
 \delta^{(3)}(\mathbf x-\mathbf y).
\end{align*}
另一方面，$\hat\pi=c^{-1}\partial_{x^0}\hat\varphi$，所以等时正则关系
$[\hat\varphi(\mathbf x),\hat\pi(\mathbf y)]=\ii\hbar\delta^{(3)}(\mathbf x-\mathbf y)$
给出同一个导数
\[
 \left.\partial_{x^0}C(x-y)\right|_{x^0=y^0}
 =-\ii\hbar c\,\delta^{(3)}(\mathbf x-\mathbf y).
\]
比较 $\delta^{(3)}$ 的系数得到
\[
 \int_0^\infty\dd\zeta\,\rho_{\rm KL}(\zeta)=1.
\]
再结合完整态插入中
$|\langle\Omega|\hat\varphi(0)|n\rangle|^2\ge0$，便得到正文
\eqnrefs{eq:KL-positive-normalization-r16} 的非负性与总权重归一化。
\end{proof}

第一条来自完整态矩阵元的模方，第二条由等时正则对易关系固定。为清楚展示归一化来源，正文先记录两种计算同一等时导数的目标式：谱表示给
\begin{align}
 \left.\partial_{x^0}C(x-y)\right|_{x^0=y^0}
 &=-\ii\hbar c
 \left[\int_0^\infty\dd\zeta\,\rho_{\rm KL}(\zeta)\right]
 \delta^{(3)}(\mathbf x-\mathbf y),
 \label{eq:KL-commutator-spectral-derivative-r15}
\end{align}
而 $C(x-y)\equiv\langle\Omega|[\hat\varphi(x),\hat\varphi(y)]|\Omega\rangle$ 与正则对易关系给
\begin{equation}
 \left.\partial_{x^0}C(x-y)\right|_{x^0=y^0}
 =-\ii\hbar c\,\delta^{(3)}(\mathbf x-\mathbf y).
 \label{eq:KL-commutator-canonical-r15}
\end{equation}
比较两式即固定\eqnrefs{eq:KL-positive-normalization-r16} 的总谱权重；详细质量壳代数放在节末推导区。

时间排序把正、负时间的 Wightman 函数拼成 Feynman 二点函数。其约化动量空间形式是
\begin{equation}
 \boxed{
 G_{F,\rm red}(p)
 =\int_0^\infty\dd\zeta\,
 \rho_{\rm KL}(\zeta)
 \frac{\ii}{p^2-\zeta+\ii0^+}.}
 \label{eq:KL-SI-r12}
\end{equation}

\begin{proposition}[二点函数的正谱表示]
设标量场的两点函数满足正性、平移不变性与正能谱条件，时间序列与质量壳测度按正文定义。则有以下结论（原式 \eqnrefs{eq:KL-SI-r12}）：
\begin{equation*}
 \boxed{
 G_{F,\rm red}(p)
 =\int_0^\infty\dd\zeta\,
 \rho_{\rm KL}(\zeta)
 \frac{\ii}{p^2-\zeta+\ii0^+}.}
\end{equation*}
\end{proposition}
\begin{proof}
固定谱参数 $\zeta$，定义
\[
 \varepsilon_{\mathbf p,\zeta}
 \equiv\sqrt{\mathbf p^2+\zeta}>0.
\]
正文\eqnrefs{eq:KL-Wightman-mass-shell-r15} 中对应这一 $\zeta$ 的正频质量壳核含时间因子
$\exp[-\ii\varepsilon_{\mathbf p,\zeta}x^0/\hbar]$。时间排序要求
\begin{align*}
 \Delta_F^{(\zeta)}(x)
 ={}&\Theta(x^0)\Delta_+^{(\zeta)}(x)
 +\Theta(-x^0)\Delta_+^{(\zeta)}(-x).
\end{align*}
关键问题是：什么四维动量分母的逆 Fourier 变换恰好产生这两个时间支路？对固定 $\mathbf p$ 考察
\begin{equation*}
 \frac{\ii}{(p^0)^2-\varepsilon_{\mathbf p,\zeta}^2+\ii0^+}.
\end{equation*}
其分布意义下的部分分式为
\begin{equation*}
 \frac{\ii}{(p^0)^2-\varepsilon^2+\ii0^+}
 =\frac{\ii}{2\varepsilon}
 \left[
 \frac1{p^0-\varepsilon+\ii0^+}
 -\frac1{p^0+\varepsilon-\ii0^+}
 \right].
\end{equation*}
因此正能极点位于 $p^0=+\varepsilon-\ii0^+$，负能极点位于
$p^0=-\varepsilon+\ii0^+$。逆 Fourier 因子是
$\exp[-\ii p^0x^0/\hbar]$：
当 $x^0>0$ 时，为使大半圆上的指数衰减，轮廓闭合到下半平面，只包围正能极点；当 $x^0<0$ 时闭合到上半平面，只包围负能极点。两次留数分别产生
\begin{align*}
 x^0>0:&\quad
 \frac{1}{2\varepsilon}
 \ee^{-\ii\varepsilon x^0/\hbar},\\
 x^0<0:&\quad
 \frac{1}{2\varepsilon}
 \ee^{+\ii\varepsilon x^0/\hbar},
\end{align*}
正好是 $\Delta_+^{(\zeta)}(x)$ 与
$\Delta_+^{(\zeta)}(-x)$ 的时间部分。把空间 Fourier 因子恢复后，固定 $\zeta$ 的时间排序核因此具有约化动量表示
\[
 \Delta_{F,\rm red}^{(\zeta)}(p)
 =\frac{\ii}{p^2-\zeta+\ii0^+}.
\]

K\"all\'en--Lehmann 表示只是把所有允许的不变量 $\zeta$ 按正谱权重线性叠加，所以
\begin{align*}
 G_{F,\rm red}(p)
 &=\int_0^\infty\dd\zeta\,
 \rho_{\rm KL}(\zeta)\Delta_{F,\rm red}^{(\zeta)}(p)\\
 &=\int_0^\infty\dd\zeta\,
 \rho_{\rm KL}(\zeta)
 \frac{\ii}{p^2-\zeta+\ii0^+},
\end{align*}
即正文\eqnrefs{eq:KL-SI-r12}。这里 $+\ii0^+$ 的作用已经由轮廓选择明确：它把正能极点放到实轴下方、负能极点放到实轴上方，从而编码“正频向未来传播、负频向过去传播”的时间排序边界条件。
\end{proof}

因此每一个精确中间态质量都对应一个自由型 Feynman 分母，而相互作用的信息全部进入正谱密度。若场与稳定单粒子态有非零重叠，谱密度包含孤立 $\delta$ 峰：
\begin{equation}
 \boxed{
 \rho_{\rm KL}(\zeta)
 =Z\,\delta(\zeta-m_{\rm phys}^2c^2)
 +\rho_{\rm KL,cont}(\zeta).}
 \label{eq:rho-pole-cont-SI-r12}
\end{equation}
其中 $Z\ge0$ 是场对该稳定单粒子态的谱权重，$m_{\rm phys}$ 是物理极点质量，$\rho_{\rm KL,cont}$ 是从最低多粒子阈值开始的连续谱。代回谱表示得到
\begin{equation}
 \boxed{
 G_{F,\rm red}(p)
 =\frac{\ii Z}{p^2-m_{\rm phys}^2c^2+\ii0^+}
 +\int_{\zeta_{\rm th}}^\infty\dd\zeta\,
 \rho_{\rm KL,cont}(\zeta)
 \frac{\ii}{p^2-\zeta+\ii0^+}.}
 \label{eq:KL-pole-plus-cont-SI-r12}
\end{equation}
若最轻连续谱通道由两个质量为 $m$ 的粒子组成，则 $\zeta_{\rm th}=4m^2c^2$。这个阈值并非只是一条抽象的“支割起点”：四维二体相空间本身已经带有无量纲开阈因子
\begin{equation}
 \boxed{
 \beta_2(\zeta)
 =\sqrt{1-\frac{4m^2c^2}{\zeta}}\,
 \Theta(\zeta-4m^2c^2).}
 \label{eq:two-particle-threshold-factor-dp68}
\end{equation}
具体的 $\rho_{\rm KL,cont}$ 还要乘上真空到该多粒子态的矩阵元模方和相应态密度，但无论相互作用细节怎样，二粒子通道都不能在 $\zeta<4m^2c^2$ 出现。\figref{fig:qft-two-particle-threshold-dp68} 直接画出\eqnrefs{eq:two-particle-threshold-factor-dp68} 的开阈行为：连续谱不是在阈值处突然获得有限常数权重，而是先由相空间从零开启。

\begin{figure}[htbp]
 \centering
 \includegraphics[width=0.72\textwidth]{qft_two_particle_threshold_dp68.pdf}
 \caption{两个等质量粒子连续谱的普适相空间开阈因子。横轴 $x=\zeta/(m^2c^2)$；虚线 $x=4$ 是二体阈值。实际 Källén--Lehmann 连续谱还包含相互作用矩阵元权重，因此本图只展示由运动学固定的阈值部分。}
 \label{fig:qft-two-particle-threshold-dp68}
\end{figure}

孤立极点与支割线因此不是两个额外模型，而是同一个精确二点函数的离散谱和连续谱贡献；LSZ 正是利用孤立极点把场插入约化为渐近单粒子态。

同一解析思想也控制线性响应。对任意厄米玻色可观测量 $\hat A$，推迟响应函数定义为
\begin{equation}
 \boxed{
 \chi_R(t)
 =-\frac{\ii}{\hbar}\Theta(t)
 \langle[\hat A(t),\hat A(0)]\rangle.}
 \label{eq:retarded-susceptibility-def-r68}
\end{equation}
因果支撑 $\chi_R(t<0)=0$ 使其 Fourier 变换在复频率上半平面解析。只要高频衰减足以使闭合轮廓积分成立，实轴边界值便满足
\begin{equation}
 \boxed{
 \begin{aligned}
 \operatorname{Re}\chi_R(\omega)
 &=\frac1\pi\mathcal P\int_{-\infty}^{\infty}\dd\omega'
 \frac{\operatorname{Im}\chi_R(\omega')}{\omega'-\omega},\\
 \operatorname{Im}\chi_R(\omega)
 &=-\frac1\pi\mathcal P\int_{-\infty}^{\infty}\dd\omega'
 \frac{\operatorname{Re}\chi_R(\omega')}{\omega'-\omega}.
 \end{aligned}}
 \label{eq:kramers-kronig-r68}
\end{equation}

\begin{proposition}[因果响应的 Hilbert 变换关系]
设响应因果且其 Fourier--Laplace 变换在上半平面解析，大圆贡献消失；若高频增长则先作足够减除。则有以下结论（原式 \eqnrefs{eq:kramers-kronig-r68}）：
\begin{equation*}
 \boxed{
 \begin{aligned}
 \operatorname{Re}\chi_R(\omega)
 &=\frac1\pi\mathcal P\int_{-\infty}^{\infty}\dd\omega'
 \frac{\operatorname{Im}\chi_R(\omega')}{\omega'-\omega},\\
 \operatorname{Im}\chi_R(\omega)
 &=-\frac1\pi\mathcal P\int_{-\infty}^{\infty}\dd\omega'
 \frac{\operatorname{Re}\chi_R(\omega')}{\omega'-\omega}.
 \end{aligned}}
\end{equation*}
\end{proposition}
\begin{proof}
由正文\eqnrefs{eq:retarded-susceptibility-def-r68} 的因果支撑 $\chi_R(t)=0$（$t<0$），其 Fourier 变换

\begin{equation*}
 \chi_R(\omega)=\int_{-\infty}^{\infty}\dd t\,\ee^{+\ii\omega t}\chi_R(t)
\end{equation*}
在复频率上半平面解析：若 $\operatorname{Im}\omega>0$，因子 $\ee^{\ii\omega t}=\ee^{\ii\operatorname{Re}\omega t}\ee^{-\operatorname{Im}\omega t}$ 保证 $t>0$ 积分收敛。对上半平面闭合积分轮廓应用 Cauchy formula，取实轴边界值后得到
\begin{equation*}
 \chi_R(\omega)
 =\frac1{\pi\ii}
 \mathcal P\int_{-\infty}^{\infty}\dd\omega'
 \frac{\chi_R(\omega')}{\omega'-\omega}.
\end{equation*}
分别取实部与虚部：
\begin{align*}
 {
 \operatorname{Re}\chi_R(\omega)
 =\frac1\pi\mathcal P\int_{-\infty}^{\infty}\dd\omega'
 \frac{\operatorname{Im}\chi_R(\omega')}{\omega'-\omega}},\\
 {
 \operatorname{Im}\chi_R(\omega)
 =-\frac1\pi\mathcal P\int_{-\infty}^{\infty}\dd\omega'
 \frac{\operatorname{Re}\chi_R(\omega')}{\omega'-\omega}}.
 \end{align*}
两式正是正文\eqnrefs{eq:kramers-kronig-r68}。
\end{proof}

这就是 Kramers--Kronig 关系：吸收与色散并不是两个独立经验函数，而是同一个因果响应的虚部和实部。

稳定粒子的极点位于物理支的实质量壳。可衰变激发与连续谱混合后，以能量 $E$ 为变量的推迟预解式在共振附近可先写成精确的局部分解
\begin{equation*}
 G_R(E)
 =\frac{Z_R}{E-E_R+\ii\Gamma_E/2}+G_{\rm reg}(E),
\end{equation*}
其中 $G_{\rm reg}$ 在所考察共振邻域解析。令 $\Delta_R$ 表示从 $E_R$ 到最近其他极点、阈值支点或背景显著变化尺度的能量距离，并定义
\begin{equation*}
 \epsilon_R
 =\frac{\max(|E-E_R|,\Gamma_E)}{\Delta_R},
 \qquad
 \epsilon_{\rm bg}
 =\frac{|E-E_R+\ii\Gamma_E/2|\,|G_{\rm reg}(E_R)|}{|Z_R|}.
\end{equation*}
当 $\epsilon_R\ll1$ 且 $\epsilon_{\rm bg}\ll1$ 时，解析背景的变化与幅度均可作为次阶修正，于是得到单极点近似
\begin{equation}
 \boxed{
 G_R(E)
 =\frac{Z_R}{E-E_R+\ii\Gamma_E/2}
 \left[1+O(\epsilon_R,\epsilon_{\rm bg})\right].}
 \label{eq:resonance-pole-SI-r12}
\end{equation}
$E_R$ 与 $\Gamma_E$ 均以 J 为单位；极点项在时域给出振幅 $\exp[-\Gamma_E t/(2\hbar)]$，因此概率衰减寿命为 $\tau=\hbar/\Gamma_E$。在相应黎曼支上解析延拓后，共振极点位于非物理支。纳米光子学准正常模、极化激元共振以及电子自能引起的能移与线宽都可用这一受控极点展开描述；当共振接近另一极点或连续谱阈值、使 $\epsilon_R$ 不再小，单极点形式必须恢复为多极点或极点--支割线的完整谱表示。

\begin{exbox}[由共振极点直接读出线宽与寿命]
在 $\epsilon_R,\epsilon_{\rm bg}\ll1$ 的单极点区间内，若探测链在线性区把响应强度乘以能量无关增益 $C_{\mathcal I}>0$，定义 $\mathcal I(E)\equiv C_{\mathcal I}|G_R(E)|^2$，则
\begin{equation}
 \mathcal I(E)
 =C_{\mathcal I}|G_R(E)|^2
 \simeq
 \frac{C_{\mathcal I}|Z_R|^2}{(E-E_R)^2+(\Gamma_E/2)^2}.
 \label{eq:resonance-lorentzian-example-dp48}
\end{equation}
峰值位于 $E=E_R$。要求信号降到峰值的一半，有
\begin{align*}
 (E-E_R)^2+(\Gamma_E/2)^2
 &=2(\Gamma_E/2)^2,\\
 |E-E_R|&=\Gamma_E/2.
\end{align*}
因此两个半高点相距
\begin{equation*}
 \Delta E_{\rm FWHM}=\Gamma_E.
\end{equation*}
另一方面，同一极点在时域给出归一化概率
$P(t)=P(0)\ee^{-\Gamma_E t/\hbar}$，故寿命
\begin{equation*}
 \tau=\frac{\hbar}{\Gamma_E}.
\end{equation*}
这说明频域线宽与时域寿命不是两个独立经验量，而是同一个复极点的实频与时间表示。对窄带光学或电子--光子共振，还可定义无量纲品质因子 $Q_R\equiv E_R/\Gamma_E$；只有在单极点条件成立时，$Q_R$ 才能由这一 Lorentz 型线宽直接解释。
\end{exbox}

% -----------------------------------------------------------------------------
% -----------------------------------------------------------------------------

\begin{exbox}[Dyson 自能与寿命]
精确二点函数的谱表示告诉我们，稳定单粒子对应实轴上的孤立极点；一旦该激发能向其它在壳通道衰变，自能获得虚部，极点离开实轴。相应的可计算对象是能量变量 $E$ 下的推迟 Green 函数
\begin{equation*}
 G^R(E)=\frac{1}{E-E_0-\Sigma^R(E)}.
\end{equation*}
在一个窄共振附近，把自能在 $E=E_r$ 邻域展开成
\begin{equation*}
 \Sigma^R(E)
 \simeq \Delta E
 +(1-Z^{-1})(E-E_r)
 -\frac{\ii\Gamma}{2Z},
\end{equation*}
其中 $\Gamma>0$ 具有能量量纲，$Z$ 是极点留数。把实部吸收到重整化共振能量 $E_r=E_0+\Delta E$ 后，Green 函数约化为
\begin{equation}
 \boxed{
 G^R(E)\simeq
 \frac{Z}{E-E_r+\ii\Gamma/2}.}
 \label{eq:dyson-pole-breit-wigner-dp65}
\end{equation}
相应谱函数定义为 $A(E)=-2\operatorname{Im}G^R(E)$，因此
\begin{equation}
 \boxed{
 A(E)=
 \frac{Z\Gamma}{(E-E_r)^2+(\Gamma/2)^2}.}
 \label{eq:dyson-lorentzian-spectral-dp65}
\end{equation}
若该窄共振几乎耗尽单粒子谱权重，则 $Z\simeq1$，并且
\begin{equation*}
 \int_{-\infty}^{\infty}\frac{\dd E}{2\pi}A(E)=Z.
\end{equation*}
所以 $\Gamma$ 不是随意画在谱线上的“展宽参数”，而是由同一个复极点的虚部固定：$\Gamma=-2\,\operatorname{Im}E_\star$，同时也是该 Lorentzian 的半峰全宽。

再把\eqnrefs{eq:dyson-pole-breit-wigner-dp65} Fourier 变回时间。对 $t>0$ 闭合下半平面，只拾取
\begin{equation*}
 E_\star=E_r-\frac{\ii\Gamma}{2}
\end{equation*}
这个极点，得到
\begin{equation*}
 G^R(t)\simeq
 -\frac{\ii Z}{\hbar}\,\theta(t)
 \exp\!\left(-\frac{\ii E_rt}{\hbar}\right)
 \exp\!\left(-\frac{\Gamma t}{2\hbar}\right),
\end{equation*}
因此振幅寿命尺度是 $2\hbar/\Gamma$，而概率或能量按
\begin{equation*}
 \ee^{-\Gamma t/\hbar}
\end{equation*}
衰减，对应通常定义的寿命
\begin{equation}
 \boxed{\tau=\frac{\hbar}{\Gamma}.}
 \label{eq:width-lifetime-si-dp65}
\end{equation}

作为数量级估算，若一个光学共振位于 $E_r=1.55\,\mathrm{eV}$，品质因子 $Q=100$，并按 $Q=E_r/\Gamma$ 定义能量线宽，则
\begin{align*}
 \Gamma&=15.5\,\mathrm{meV},\\
 \tau&=\frac{\hbar}{\Gamma}
 \simeq42.5\,\mathrm{fs}.
\end{align*}
“复极点 $\leftrightarrow$ Lorentzian 频谱 $\leftrightarrow$ 指数时间衰减”是推迟预解式的一般解析结构；有损光学正常模、QNM Green 张量和电子能量损失谱都由相应自能的复极点实现同一结构。


\medskip

稳定粒子对应实轴上的孤立极点；当该态能够衰变时，推迟 Green 函数的极点移入下半复能量平面。\eqnrefs{eq:dyson-pole-breit-wigner-dp65} 已把窄共振附近的分母写成 Breit--Wigner 形式。对能量变量作 Fourier 变换时，下半平面的极点直接产生时间域指数衰减，并给出\eqnrefs{eq:width-lifetime-si-dp65} 的 $\tau=\hbar/\Gamma$。这一近似要求共振孤立、$\Gamma\ll E_r$，且自能在宽度 $\Gamma$ 的能量范围内变化缓慢。


从 Dyson 分母
\begin{equation*}
 D(E)=E-E_0-\Sigma^R(E)
\end{equation*}
开始。把自能实部在 $E_r$ 附近展开：
\begin{equation*}
 \operatorname{Re}\Sigma^R(E)
 \simeq \operatorname{Re}\Sigma^R(E_r)
 +(E-E_r)\,\partial_E\operatorname{Re}\Sigma^R(E_r).
\end{equation*}
令
\begin{equation*}
 E_r-E_0-\operatorname{Re}\Sigma^R(E_r)=0,
 \qquad
 Z^{-1}=1-\partial_E\operatorname{Re}\Sigma^R(E_r),
\end{equation*}
并定义
\begin{equation*}
 \Gamma=-2Z\operatorname{Im}\Sigma^R(E_r)>0.
\end{equation*}
于是
\begin{align*}
 D(E)
 &\simeq Z^{-1}(E-E_r)+\ii\frac{\Gamma}{2Z},\\
 G^R(E)=\frac1{D(E)}
 &\simeq\frac{Z}{E-E_r+\ii\Gamma/2},
\end{align*}
即\eqnrefs{eq:dyson-pole-breit-wigner-dp65}。

取虚部：
\begin{align*}
 \operatorname{Im}\frac{1}{x+\ii a}
 &=-\frac{a}{x^2+a^2},\\
 A(E)=-2\operatorname{Im}G^R(E)
 &=\frac{Z\Gamma}{(E-E_r)^2+(\Gamma/2)^2},
\end{align*}
得到\eqnrefs{eq:dyson-lorentzian-spectral-dp65}。再计算时间核
\begin{equation*}
 G^R(t)=\int_{-\infty}^{\infty}\frac{\dd E}{2\pi\hbar}
 \ee^{-\ii Et/\hbar}G^R(E).
\end{equation*}
对 $t>0$，因 $\ee^{-\ii Et/\hbar}$ 在下半平面衰减，闭合下半圆；唯一极点是 $E_\star=E_r-\ii\Gamma/2$。留数定理给
\begin{equation*}
 G^R(t)\simeq
 -\frac{\ii Z}{\hbar}\,
 \ee^{-\ii E_rt/\hbar}\ee^{-\Gamma t/(2\hbar)}.
\end{equation*}
振幅模方因而为
\begin{equation*}
 |G^R(t)|^2\simeq
 \frac{|Z|^2}{\hbar^2}\ee^{-\Gamma t/\hbar}.
\end{equation*}
若把寿命定义为概率下降到 $\ee^{-1}$ 的时间，便得到\eqnrefs{eq:width-lifetime-si-dp65}。
\end{exbox}






\begin{derivation}[全入射振幅与物理散射通道]
全入射约定把每一条外腿的四动量都定义成指向相互作用区。设全入射核的四个参数为 $k_i$，则正文\eqnrefs{eq:all-incoming-amplitude-dp9} 要求
\[
k_1+k_2+k_3+k_4=0.
\]
对物理过程 $1+2\to3+4$，实验四动量满足
\[
p_1+p_2=p_3+p_4.
\]
前两条腿本来就流入，所以
\[
k_1=p_1,\qquad k_2=p_2.
\]
末态 3、4 的物理箭头流出相互作用区；若仍要求 $k_i$ 全部向内，必须反号：
\[
k_3=-p_3,\qquad k_4=-p_4.
\]
于是全入射的动量守恒自动变成
\begin{align*}
k_1+k_2+k_3+k_4
&=p_1+p_2-p_3-p_4\\
&=0,
\end{align*}
与物理守恒完全一致。将四个 $k_i$ 代回同一个截肢四点核，便得到正文\eqnrefs{eq:physical-from-all-incoming-dp9}：
\[
\mathcal M_{12\to34}
=\mathcal A(p_1,p_2,-p_3,-p_4).
\]
\end{derivation}

\begin{derivation}[低速散射振幅与势能的匹配]
低速弹性区有 $|\vb{p}_f|=|\vb{p}_i|$。定义非相对论振幅为阈值极限
\begin{equation*}
\mathcal M_{\rm rel}^{\rm NR}(\vb{q})
\equiv\lim_{|\vb{p}_i|/(m_{A,B}c)\to0}\mathcal M_{\rm rel},
\qquad
\lim_{|\vb{p}_i|\to0}\sqrt{s}=(m_A+m_B)c.
\end{equation*}
因此对截面主式取同一极限得到
\begin{equation*}
 \frac{\dd\sigma}{\dd\Omega}
 =\frac{\hbar^2}{64\pi^2(m_A+m_B)^2c^2}
 |\mathcal M_{\rm rel}^{\rm NR}(\vb{q})|^2.
\end{equation*}
另一方面，对\eqnrefs{eq:sch-lippmann-r13} 作一次 Born 迭代，采用同一出射波相位约定可写
\begin{equation*}
 f_B(\vb{q})
 =-\frac{m_{\rm red}}{2\pi\hbar^2}
 \widetilde V_{\rm NR}(\vb{q}),
 \qquad
 \frac{\dd\sigma}{\dd\Omega}=|f_B|^2,
\end{equation*}
其中 $m_{\rm red}=m_Am_B/(m_A+m_B)$。与相对论截面对应的散射振幅归一化为
\begin{equation*}
 f_{\rm rel}(\vb{q})
 =\frac{\hbar}{8\pi(m_A+m_B)c}
 \mathcal M_{\rm rel}^{\rm NR}(\vb{q}),
\end{equation*}
其相位由前面固定的 $S$ 矩阵约定与出射 Green 核共同确定。令 $f_B=f_{\rm rel}$，得到
\begin{align*}
 \widetilde V_{\rm NR}(\vb{q})
 &=-\frac{\hbar^3}{4m_{\rm red}(m_A+m_B)c}
 \mathcal M_{\rm rel}^{\rm NR}(\vb{q})\\
 &=-\frac{\hbar^3}{4m_Am_Bc}
 \mathcal M_{\rm rel}^{\rm NR}(\vb{q}),
\end{align*}
即\eqnrefs{eq:relativistic-to-nr-potential-matching-dp75}。
\end{derivation}

\begin{derivation}[质心系中的二体相空间]
正文\eqnrefs{eq:phase2-si-start-r9} 仍含两个末态三动量积分和一个四维守恒 $\delta$ 函数。在质心系中，三维 $\delta$ 函数先固定 $\vb{p}_4=-\vb{p}_3$，剩下的能量 $\delta$ 函数再固定径向动量模；最后只留下末态方向 $\dd\Omega$，从而得到正文\eqnrefs{eq:phase2-final-si-r9}。

采用能量四矢量
\begin{equation*}
 \mathcal P^\mu\equiv cp^\mu=(E,c\vb{p}),
 \end{equation*}
它只是物理四动量乘以常数 $c$。定义标准二体相空间因子
\begin{equation*}
 \dd\mathfrak P_2
 =(2\pi)^4\delta^{(4)}(\mathcal P_1+\mathcal P_2-\mathcal P_3-\mathcal P_4)
 \prod_{a=3,4}\frac{c^3\dd^3 p_a}{(2\pi)^3\,2E_a}.
 \end{equation*}
每个 $c^3\dd^3 p$ 都是 $\dd^3(c\vb{p})$，所以 \eqnrefs{eq:phase2-si-start-r9} 完全由能量量纲的四动量组成。在质心系 $\mathcal P_1+\mathcal P_2=(c\sqrt s,\mathbf0)$。三维 $\delta$ 先令 $\vb{p}_4=-\vb{p}_3$，于是
\begin{equation*}
 \dd\mathfrak P_2
 =\frac{c^4p^2\dd p\,\dd\Omega}{(2\pi)^2 4E_3E_4}
 \delta(c\sqrt s-E_3-E_4).
\end{equation*}
令
\begin{equation*}
 f(p)=c\sqrt s-E_3(p)-E_4(p).
\end{equation*}
由 $\dd E_a/\dd p=c^2p/E_a$，
\begin{equation*}
 \left|f'(p)\right|
 =c^2p\left(\frac1{E_3}+\frac1{E_4}\right)
 =\frac{c^3p\sqrt s}{E_3E_4}.
\end{equation*}
所以
\begin{equation*}
 \delta[f(p)]
 =\frac{E_3E_4}{c^3p_f\sqrt s}\delta(p-p_f).
\end{equation*}
代回并积分 $p$，得到
\begin{equation*}
 {
 \dd\mathfrak P_2
 =\frac1{16\pi^2}\frac{|\vb{p}_f|}{\sqrt s}\,\dd\Omega.}
 \end{equation*}
右端无量纲，因为 $|\vb{p}_f|$ 与 $\sqrt s$ 都具有动量量纲。
\end{derivation}

\begin{derivation}[不变入射通量与相遇速率]
对两条入射轨迹，普通同一参考系中的相对相遇速率不能简单写成任意参考系都适用的 $|\vb{v}_1-\vb{v}_2|$。由四动量构造 Lorentz 标量

\begin{equation*}
 \mathcal I_{12}
 \equiv\sqrt{(p_1\!\cdot p_2)^2-m_1^2m_2^2c^4}.
 \end{equation*}
它具有动量平方量纲。定义 M\o ller 相对速率
\begin{equation*}
 v_{\rm M}
 \equiv
 \frac{c^3\mathcal I_{12}}{E_1E_2}.
 \end{equation*}
检查质心系时不跳过中间平方。写
\begin{equation*}
 p_1^\mu=\left(\frac{E_1}{c},\vb{p}_i\right),
 \qquad
 p_2^\mu=\left(\frac{E_2}{c},-\vb{p}_i\right),
\end{equation*}
于是
\begin{equation*}
 p_1\!\cdot p_2=\frac{E_1E_2}{c^2}+|\vb{p}_i|^2.
\end{equation*}
两条质量壳又分别给
\begin{equation*}
 m_a^2c^2=\frac{E_a^2}{c^2}-|\vb{p}_i|^2,
 \qquad a=1,2.
\end{equation*}
因此 \eqnrefs{eq:invariant-flux-scalar-r9} 根号中的量逐项化为
\begin{align*}
 \mathcal I_{12}^2
 &=\left(\frac{E_1E_2}{c^2}+p_i^2\right)^2
 -\left(\frac{E_1^2}{c^2}-p_i^2\right)
  \left(\frac{E_2^2}{c^2}-p_i^2\right),\\
 &=\frac{E_1^2E_2^2}{c^4}
 +\frac{2E_1E_2p_i^2}{c^2}+p_i^4
 -\frac{E_1^2E_2^2}{c^4}
 +\frac{(E_1^2+E_2^2)p_i^2}{c^2}-p_i^4,\\
 &=\frac{p_i^2}{c^2}(E_1+E_2)^2,
\end{align*}
其中 $p_i\equiv|\vb{p}_i|\ge0$。取正根得到
\begin{equation*}
 {
 \mathcal I_{12}
 =\frac{|\vb{p}_i|(E_1+E_2)}{c}
 =|\vb{p}_i|\sqrt s.}
 \end{equation*}
最后一步用了 $E_1+E_2=c\sqrt s$。于是
\begin{align*}
 v_{\rm M}
 &=|\vb{p}_i|c^2
 \left(\frac1{E_1}+\frac1{E_2}\right),\\
 &=|\vb{v}_1-\vb{v}_2|,
\end{align*}
因为 $\vb{v}_a=c^2\vb{p}_a/E_a$。所以 \eqnrefs{eq:moller-velocity-si-r9} 的协变表达确实回到熟悉的相对速度。
\end{derivation}

\begin{derivation}[盒归一化怎样给出二体散射截面]

取空间周期盒体积 $V$ 与观测时间 $T$。空间连续态满足
\begin{equation*}
 (2\pi\hbar)^3\delta^{(3)}(\mathbf0)=V,
 \qquad
 \sum_{\vb{p}}\longrightarrow
 V\int\frac{\dd^3 p}{(2\pi\hbar)^3}.
\end{equation*}
时间方向令 $\delta E\equiv E_f-E_i$，则
\begin{equation*}
 \int_{-T/2}^{T/2}\dd t\,\ee^{\ii\delta E t/\hbar}
 =T\,\operatorname{sinc}\!\left(\frac{\delta E T}{2\hbar}\right),
\end{equation*}
并由正文\eqnrefs{eq:delta-square-time-r15} 有
\begin{equation*}
 \lim_{T\to\infty}\frac1T
 \left|\int_{-T/2}^{T/2}\dd t\,\ee^{\ii\delta E t/\hbar}\right|^2
 =2\pi\hbar\delta(\delta E).
\end{equation*}
因此四维动量守恒 $\delta$ 函数平方在概率除以 $T$ 后留下一个能量守恒 $\delta$，其空间零点因子给出一个 $V$。结合两条初态的协变归一化，单位盒中单对入射粒子的二体末态跃迁率为正文\eqnrefs{eq:box-transition-rate-r15}，即
\begin{equation*}
 \dd\Gamma_{1+2\to2}
 =\frac{\hbar^2c^3}{4VE_1E_2}
 \overline{|\mathcal M|^2}\,\dd\mathfrak P_2.
\end{equation*}

单个归一化入射粒子在盒中的数密度是 $1/V$。正文\eqnrefs{eq:moller-velocity-si-r9} 给出相对速率，所以入射通量密度为
\begin{equation*}
 j_{12}=\frac{v_{\rm M}}{V}
 =\frac{c^3\mathcal I_{12}}{VE_1E_2}.
\end{equation*}
按正文中“跃迁率除以入射通量密度”的截面关系：
\begin{align*}
 \dd\sigma
 &=\frac{\dd\Gamma_{1+2\to2}}{j_{12}}\\
 &=\frac{\hbar^2}{4\mathcal I_{12}}
 \overline{|\mathcal M|^2}\,\dd\mathfrak P_2.
\end{align*}
$V$、$T$ 与参考系依赖的 $E_1E_2$ 在这一步全部消去。对 $2\to2$ 过程代入正文\eqnrefs{eq:phase2-final-si-r9}，并在质心系使用
$\mathcal I_{12}=|\vb{p}_i|\sqrt s$，得到
\begin{align*}
 \frac{\dd\sigma}{\dd\Omega}
 &=\frac{\hbar^2}{4|\vb{p}_i|\sqrt s}
 \overline{|\mathcal M|^2}
 \frac1{16\pi^2}\frac{|\vb{p}_f|}{\sqrt s}\\
 &=\frac{\hbar^2}{64\pi^2s}
 \frac{|\vb{p}_f|}{|\vb{p}_i|}
 \overline{|\mathcal M|^2},
\end{align*}
即正文\eqnrefs{eq:2to2-cross-section-master-si-r9}。
\end{derivation}

\begin{derivation}[单粒子衰变率的归一化]
协变态归一化的单粒子初态在盒归一化后带因子 $(2P^0V)^{-1/2}$； 末态每条外腿给对应的质量壳测度。 LSZ 矩阵元中的四维 $\delta$ 在平方后，空间部分利用

\begin{equation*}
 (2\pi\hbar)^3\delta^{(3)}(\mathbf0)=V
\end{equation*}
消去母粒子的盒体积，而时间部分利用 \eqnrefs{eq:delta-square-time-r15} 把 $|\delta(E_f-E_i)|^2/T$ 化为单个能量 $\delta$。于是概率除以 $T$ 后只剩一个母粒子归一化因子 $1/(2P^0)$ 与 $n$ 体末态相空间。

为避免混淆能量量纲和动量量纲，先定义能量归一化截肢振幅
\begin{equation*}
 \mathcal A_E\equiv c\,\mathcal M.
\end{equation*}
若 $\mathcal M$ 在某个过程具有一个动量量纲，则 $\mathcal A_E$ 具有能量量纲。有限盒推导给
\begin{equation*}
 \dd\Gamma_{1\to n}
 =\frac{1}{2E_P\hbar}\,
 \overline{|\mathcal A_E|^2}\,\dd\mathfrak P_n,
 \qquad E_P=Mc^2.
 \end{equation*}
把 $\mathcal A_E=c\mathcal M$ 与 $E_P=Mc^2$ 代入，两个 $c^2$ 消去：
\begin{align*}
 \dd\Gamma_{1\to n}
 &=\frac{c^2}{2Mc^2\hbar}
 \overline{|\mathcal M|^2}\,\dd\mathfrak P_n,\\
 &={
 \frac{1}{2M\hbar}
 \overline{|\mathcal M|^2}\,\dd\mathfrak P_n.}
 \end{align*}
这里 $\dd\mathfrak P_n$ 使用能量四矢量 $\mathcal P^\mu=cp^\mu$ 的定义。以无量纲 Yukawa 耦合的费米子对衰变为例，振幅含有双线性因子 $\bar u v$；在上述外腿归一化下该双线性具有动量量纲，因此
\begin{equation*}
 \left[\frac{|\mathcal M|^2}{M\hbar}\right]
 =\mathrm{s^{-1}},
\end{equation*}
与衰变率的 SI 单位一致。这样 \eqnrefs{eq:master-decay-rate-si-r9} 与散射截面主式共享同一套有限盒、有限时间和外腿归一化，而不是独立记忆的另一条规则。
\end{derivation}

\begin{derivation}[三体相空间与 Dalitz 变量]
在母粒子静止系中，对正文\eqnrefs{eq:phase3-recursive-dp52} 的两个二体因子分别使用\eqnrefs{eq:phase2-final-si-r9}：

\begin{equation*}
\dd\mathfrak P_3
=\frac{\dd\mathcal S_{12}}{2\pi}
\frac{1}{16\pi^2}\frac{|\vb{p}_3|}{Mc}\dd\Omega_3
\frac{1}{16\pi^2}\frac{|\vb{p}_2^*|}{\sqrt{\mathcal S_{12}}/c}\dd\Omega_2^*.
\end{equation*}
若母粒子无自旋或已做自旋平均，整体方向积分给 $\int\dd\Omega_3=4\pi$。在 $q_{12}$ 静止系中以粒子 3 为极轴，方位角再给 $2\pi$，于是
\begin{equation*}
\dd\mathfrak P_3
=\frac{\dd\mathcal S_{12}\,\dd\cos\theta^*}{64\pi^3}
\frac{|\vb{p}_3|}{Mc}
\frac{|\vb{p}_2^*|}{\sqrt{\mathcal S_{12}}/c}.
\end{equation*}
正文\eqnrefs{eq:dalitz-angular-map-dp68} 给出
\begin{equation*}
\left|\frac{\partial\mathcal S_{23}}{\partial\cos\theta^*}\right|
=2c^2|\vb{p}_2^*||\vb{p}_3^*|.
\end{equation*}
因此把 $\cos\theta^*$ 换成 $\mathcal S_{23}$ 时，两个动量模恰好与 Jacobian 相消；再用 $\sqrt{\mathcal S_{12}}$ 和母粒子质量整理剩余因子，得到
\begin{equation*}
\dd\mathfrak P_3
=\frac{\dd\mathcal S_{12}\,\dd\mathcal S_{23}}{128\pi^3M^2c^4},
\end{equation*}
即正文\eqnrefs{eq:phase3-dalitz-dp52}。
\end{derivation}

\begin{derivation}[二体阈值附近的相空间因子]

正文\eqnrefs{eq:bubble-imag-phase-space-dp8} 表明虚部只有在两个内部粒子同时满足正能质量壳和能量守恒时才存在。对等质量情形进入质心系，能量 $\delta$ 函数固定唯一的径向动量 $k_*$；解出这个根并把它写成相对论速度因子，就得到正文\eqnrefs{eq:two-body-threshold-beta-dp8}。

在质心系 $\vb{P}=0$ 且 $m_1=m_2=m$ 时，
\begin{equation*}
 E_1=E_2=E_k=\sqrt{m^2c^4+c^2k^2},
 \qquad
 \sqrt{s}\,c=2E_k.
\end{equation*}
平方后有
\begin{align*}
 sc^2
 &=4m^2c^4+4c^2k_*^2,\\
 k_*^2
 &=\frac{s-4m^2c^2}{4}.
\end{align*}
因此只有 $s\ge4m^2c^2$ 时 $k_*$ 才为实数。再用 $E_{\rm cm}=\sqrt{s}\,c$，
\begin{align*}
 \beta_*(s)
 &=\frac{ck_*}{E_{\rm cm}/2}
 =\frac{2ck_*}{\sqrt{s}\,c}\\
 &=\sqrt{1-\frac{4m^2c^2}{s}},
\end{align*}
即正文\eqnrefs{eq:two-body-threshold-beta-dp8}。为把阈值因子的整体归一化也固定下来，定义正文\eqnrefs{eq:bubble-imag-phase-space-dp8} 中的径向相空间核
\begin{equation*}
 \rho_2(s)
 \equiv
 \int\frac{\dd^3\vb{k}}{(2\pi\hbar)^3}
 \frac{c^2}{4E_k^2}
 \delta\!\left(\sqrt{s}\,c-2E_k\right).
\end{equation*}
利用
\begin{equation*}
 \left.\frac{\dd(2E_k)}{\dd k}\right|_{k_*}
 =\frac{2c^2k_*}{E_*},
 \qquad E_*\equiv E_{k_*}=\frac{\sqrt{s}\,c}{2},
\end{equation*}
径向 $\delta$ 积分严格给出
\begin{align*}
 \rho_2(s)
 &=\frac{4\pi k_*^2}{(2\pi\hbar)^3}
 \frac{c^2}{4E_*^2}
 \frac{E_*}{2c^2k_*}\\
 &=\frac{k_*}{16\pi^2\hbar^3E_*}
 =\frac{\beta_*(s)}{16\pi^2\hbar^3c}.
\end{align*}
因此连续谱在阈值处按确定的线性因子 $\beta_*(s)$ 从零打开。
\end{derivation}

